\subsection{Positive linear forms}

DEFINITION 2. --- Let A be an involutive algebra. A linear form \(\lambda\) on A is said to be positive if \(\lambda(a^{*}a) \in \mathbf{R}_{+}\) for all \(a \in \mathrm{A}\) .

If A is an involutive Banach algebra, we denote by \(A_{+}^{\prime}\) the set of continuous positive linear forms on A.

Let A be an involutive Banach algebra. The set \(A_{+}^{\prime}\) is a pointed convex cone in the real vector space of C-linear forms on A.

Lemma 2. --- Let A be an involutive Banach algebra and let \(\lambda\) be a positive linear form on A.

\begin{enumerate}
\def\labelenumi{\alph{enumi})}
\tightlist
\item
  For all \(a\) and \(b\) in A, we have \(\lambda(a^{*}b)=\overline{\lambda(b^{*}a)}\) and
\end{enumerate}

\[
\left| \lambda \left(a ^ {*} b\right) \right| ^ {2} \leqslant \lambda \left(a ^ {*} a\right) \lambda \left(b ^ {*} b\right);
\]

\begin{enumerate}
\def\labelenumi{\alph{enumi})}
\setcounter{enumi}{1}
\tightlist
\item
  If A is unital, then the linear form \(\lambda\) is continuous and its norm is equal to \(\lambda(1)\) .
\end{enumerate}

The map \((a,b)\mapsto\lambda(a^{*}b)\) is a positive Hermitian form on A; it therefore satisfies \(\lambda(a^{*}b)=\overline{\lambda(b^{*}a)}\) and \(|\lambda(a^{*}b)|^{2}\leqslant\lambda(a^{*}a)\lambda(b^{*}b)\) for all a and b in A (EVT, V, p. 2, remark and p. 3, prop. 2).

Let us prove \(b\) ). Let \(a \in A\) be a Hermitian element of norm \(< 1\) . The spectral radius of \(a\) is less than \(\|a\|\) , hence the spectrum of \(a\) is contained in the open unit disk D of C (th. 1 of I, p. 24). Let \(f\) be a holomorphic function on D such that \(f(z)^{2} = 1 - z\) for all \(z \in D\) (Lemma 1 of V, p. 435). Apply the holomorphic functional calculus to the element \(a\) and to the function \(f\) . The element \(b = f(a)\) satisfies \(b^{2} = 1 - a\) (I, p. 74, Thm. 5). By Prop. 3 of V, p. 435, we also have \(f(a)^{*} = f^{*}(a^{*}) = f(a)\) , hence \(b\) is Hermitian. It then follows that \(\lambda(1 - a) = \lambda(b^{*}b) \geqslant 0\) , whence \(\lambda(a) \leqslant \lambda(1)\) .

Now let \(b \in A\) be of norm \(< 1\) . The element \(b^*b\) is Hermitian and \(\|b^*b\| < 1\) , so applying \(a\) ) with \(a = 1\) , we find

\[
| \lambda (b) | ^ {2} \leqslant \lambda (1) \lambda (b ^ {*} b) \leqslant \lambda (1) ^ {2},
\]

with equality if \(b = 1\) . The linear form \(\lambda\) is therefore continuous, and its norm is equal to \(\lambda(1)\) . The assertion \(b\) ) is proved.

Example. --- Let X be a compact topological space and let A be the star algebra \(\mathcal{C}(\mathrm{X})\) . Positive linear forms on A are identified with positive measures on X (INT, III, p. 52, § 1, n° 6, th. 1).

Lemma 3. --- Let A be a unital involutive Banach algebra admitting an approximate unit (I, p. 120, def. 7). Let λ be a continuous positive linear form on A.

\begin{enumerate}
\def\labelenumi{\alph{enumi})}
\item
  For all a in A, we have \(\lambda(a^{*})=\overline{\lambda(a)}\) and \(|\lambda(a)|^{2}\leqslant\|\lambda\|\lambda(a^{*}a)\) ;
\item
  Let \(\widetilde{\mathrm{A}}\) be the involutive algebra obtained from A by adjoining a unit element, and let e be its unit element. There exists a continuous positive linear form \(\widetilde{\lambda}\) on \(\widetilde{\mathrm{A}}\) which extends \(\lambda\) and such that \(\widetilde{\lambda}(e) = \|\lambda\|\) ;
\item
  For all \(a\) and \(b\) in A, we have \(|\lambda(b^{*}ab)|\leqslant\|a\|\lambda(b^{*}b)\) .
\end{enumerate}

Let us prove \(a\) ). Let \(\mathfrak{F}\) be an approximate unit of A. Let \(a \in \mathrm{A}\) . Using Lemma 2, \(a\) ) and the definition of an approximate unit, one finds

\[
\lambda (a ^ {*}) = \lim _ {f, \mathfrak {F}} \lambda (f a ^ {*}) = \lim _ {f, \mathfrak {F}} \overline {{\lambda (a f ^ {*})}} = \lim _ {f, \mathfrak {F}} \overline {{\lambda ((f a ^ {*}) ^ {*})}} = \overline {{\lambda (a)}},
\]

whence (loc. cit.)

\[
| \lambda (a) | ^ {2} = \lim _ {f, \mathfrak {F}} | \lambda (f a) | ^ {2} \leqslant \lambda (a ^ {*} a) \operatorname * {limsup} _ {f, \mathfrak {F}} \lambda (f f ^ {*}) \leqslant \| \lambda \| \lambda (a ^ {*} a),
\]

since \(\mathfrak{F}\) is a filter on the unit ball of A.

To prove \(b\) ), one may assume that \(\lambda\) is not zero, and then that \(\| \lambda \| = 1\) . For \(a \in \mathrm{A}\) and \(z \in \mathbf{C}\) , set \(\widetilde{\lambda}(a + z \cdot e) = \lambda(a) + z\) . The map \(\widetilde{\lambda}\) is a continuous linear form on \(\widetilde{\mathrm{A}}\) which extends \(\lambda\) and satisfies \(\widetilde{\lambda}(e) = 1\) . It is positive: for all \(a \in \mathrm{A}\) and \(z \in \mathbf{C}\) , according to \(a\) ), one computes

\[
\begin{array}{r l} \widetilde {\lambda} ((a + z \cdot e) ^ {*} (a + z \cdot e)) & = \lambda (a ^ {*} a) + \overline {{z}} \lambda (a) + z \lambda (a ^ {*}) + | z | ^ {2} \\ & = | z + \lambda (a) | ^ {2} + \lambda (a ^ {*} a) - | \lambda (a) | ^ {2} \geqslant 0. \end{array}
\]

Finally, let \(b\) be an element of A. The map \(a \mapsto \widetilde{\lambda}(b^{*}ab)\) is a positive linear form on the unital involutive Banach algebra \(\widetilde{A}\) . It is therefore continuous with norm equal to its value at \(e\) (Lemma 2, \(b\) ), which is equal to \(\lambda(b^{*}b)\) , whence assertion \(c\) ).

PROPOSITION 4. — Let A be an involutive Banach algebra.

\begin{enumerate}
\def\labelenumi{\alph{enumi})}
\item
  For every Hilbert space E, every morphism of involutive algebras \(\varphi\colon A\to\mathcal{L}(E)\) and every vector \(x\in E\) , the linear form defined on A by \(\lambda(a)=\langle x|\varphi(a)x\rangle\) is a continuous positive linear form;
\item
  Let \(\lambda \in \mathrm{A}_{+}^{\prime}\) . The map \(f\) from \(\mathrm{A} \times \mathrm{A}\) to \(\mathbf{C}\) defined by \(f(a, b) = \lambda (a^{*}b)\) for all \((a, b) \in \mathrm{A} \times \mathrm{A}\) is a universally positive kernel on \(\mathrm{A}\) .
\end{enumerate}

Let us prove \(a\) ). For every \(a \in \mathrm{A}\) , we have

\[
\lambda (a ^ {*} a) = \langle x \mid \varphi (a ^ {*} a) x \rangle = \| \varphi (a) x \| ^ {2} \geqslant 0
\]

hence \(\lambda\) is a positive linear form on A; it is continuous since \(\varphi\) is continuous (Prop. 2 of I, p. 104).

Let us prove \(b\) ). The function \(f\) is continuous. For every \(n \in \mathbf{N}\) , every family \((a_i)_{0 \leqslant i \leqslant n}\) in A and every family \((t_i)_{0 \leqslant i \leqslant n}\) of complex numbers, we have

\[
\sum_ {i = 0} ^ {n} \sum_ {j = 0} ^ {n} \bar {t} _ {i} t _ {j} f (a _ {i}, a _ {j}) = \lambda \left(\left(\sum_ {i = 0} ^ {n} t _ {i} a _ {i}\right) ^ {*} \left(\sum_ {j = 0} ^ {n} t _ {j} a _ {j}\right)\right) \geqslant 0,
\]

whence the result (Def. 1 of V, p. 433).

DEFINITION 3. — Let A be an involutive Banach algebra. Let \(\lambda \in \mathrm{A}_{+}^{\prime}\) be a continuous positive linear form on A. A Hilbertian realization of \(\lambda\) is a triplet (E, \(x, \varphi\) ) consisting of a Hilbert space E, an element \(x\) of E, and a morphism of involutive algebras \(\varphi\) from A into \(\mathcal{L}(\mathrm{E})\) such that \(\lambda(b^{*}a) = \langle \varphi(b)x | \varphi(a)x \rangle\) for all \((a,b) \in \mathrm{A}^{2}\) .

If the set of elements \(\varphi(a)x\) for \(a\in\mathrm{A}\) is total in E, we say that (E, \(x,\varphi\) ) is a cyclic Hilbert-space realization of \(\lambda\) .

PROPOSITION 5. --- Let A be an involutive Banach algebra admitting an approximate unit. Every continuous positive linear form on A admits a cyclic Hilbert realization (E, x, φ). If A is unital, then there exists such a realization where the morphism φ is unital.

We may assume that A is a unital algebra (Lemma 3, b)).

Let \(\lambda\) be a positive continuous linear form on A. Let (E, \(g\) ) be a cyclic Hilbertian realization of the universally positive kernel \(f\) on A defined by \(f(a,b) = \lambda(a^{*}b)\) (prop. 4 and th. 1). The map \(g\) is continuous, its image is total in E, and we have \(\lambda(a^{*}b) = \langle g(a)|g(b)\rangle\) for all \((a,b) \in \mathrm{A}^2\) . Set \(x = g(1) \in \mathrm{E}\) .

The map g from A into E is linear: indeed, for \((a,b,c)\in\mathrm{A}^{3}\) and \((s,t)\in\mathbf{C}^{2}\) , it follows that

\[
\begin{array}{c} \langle g (c) | g (s a + t b) \rangle = \lambda (c ^ {*} (s a + t b)) \\ = s \lambda (c ^ {*} a) + t \lambda (c ^ {*} b) = \langle g (c) | s g (a) + t g (b) \rangle , \end{array}
\]

whence the assertion since the image of \(g\) is total in E.

In particular, the image F of g is a dense vector subspace of E. The kernel of g is a left ideal of A: if \(g(b) = 0\) , then for all a and c in A, it follows that

\[
\langle g (a b) \mid g (c) \rangle = \lambda ((a b) ^ {*} c) = \lambda (b ^ {*} (a ^ {*} c)) = \langle g (b) \mid g (a ^ {*} c) \rangle = 0,
\]

hence \(g(ab) = 0\) since F is dense in E.

Let \(a \in \mathrm{A}\) . Since the kernel of \(g\) is a left ideal of \(\mathrm{A}\) , there exists a linear map \(\widetilde{\varphi}(a) \colon \mathrm{F} \to \mathrm{F}\) such that \(\widetilde{\varphi}(a)(g(b)) = g(ab)\) for all \(b \in \mathrm{A}\) . Moreover, for every \(b \in \mathrm{A}\) , it follows that

\[
\| \widetilde {\varphi} (a) (g (b)) \| ^ {2} = \| g (a b) \| ^ {2} = \lambda \left(b ^ {*} a ^ {*} a b\right) \leqslant \| a ^ {*} a \| \lambda \left(b ^ {*} b\right) \leqslant \| a \| ^ {2} \| g (b) \| ^ {2}
\]

(lemma 3, c)), hence \(\widetilde{\varphi}(a)\) is continuous and of norm \(\leqslant \|a\|\) . Consequently, there exists a unique endomorphism \(\varphi(a) \in \mathcal{L}(\mathrm{E})\) which induces \(\widetilde{\varphi}(a)\) by passage to the subspaces.

Let \(a\) and \(b\) be in A. We have

\[
(\varphi (a) \circ \varphi (b)) (g (c)) = g (a b c) = \varphi (a b) (g (c))
\]

for all c in A, hence \(\varphi(ab) = \varphi(a) \circ \varphi(b)\) since F is dense in E.

The map \(\varphi\) from A into \(\mathcal{L}(\mathrm{E})\) is linear: indeed, for every \((a,b)\in\mathrm{A}^{2}\) and every \((s,t)\in\mathbf{C}^{2}\) , we have

\[
\varphi (s a + t b) (g (c)) = g ((s a + t b) c) = (s \varphi (a) + t \varphi (b)) g (c)
\]

for all \(c \in A\) , whence \(\varphi(sa + tb) = s\varphi(a) + t\varphi(b)\) since F is dense in E. Since \(\|\varphi(b)\| \leqslant \|b\|\) for all \(b \in A\) , the map \(\varphi\) is continuous. We also have \(\varphi(1) = 1_{\mathrm{E}}\) since \(\varphi(1)(g(a)) = g(a)\) for all \(a \in A\) and since F is dense in E.

Finally, let a, b, c be in A. We have

\[
\begin{array}{c} \langle g (c) \mid \varphi (a ^ {*}) (g (b)) \rangle = \langle g (c) \mid g (a ^ {*} b) \rangle = \lambda (c ^ {*} a ^ {*} b) \\ = \langle g (a c) \mid g (b) \rangle = \langle \varphi (a) (g (c)) \mid g (b) \rangle \end{array}
\]

whence \(\varphi(a^{*})=\varphi(a)^{*}\) since F is dense in E.

In conclusion, the map \(\varphi\) is a continuous morphism of involutive algebras from A into \(\mathcal{L}(\mathrm{E})\) and \(g(a) = \varphi(a)x\) for all \(a \in \mathrm{A}\) ; consequently, the triple \((\mathrm{E}, x, \varphi)\) is a cyclic Hilbert realization of \(\lambda\) .

PROPOSITION 6. --- Let A be an involutive Banach algebra and let λ be a continuous positive linear form on A. Let (E₁, x₁, φ₁) and (E₂, x₂, φ₂) be Hilbertian realizations of λ. Suppose that (E₁, x₁, φ₁) is a cyclic Hilbertian realization.

\begin{enumerate}
\def\labelenumi{\alph{enumi})}
\item
  There exists a unique continuous linear map u from \(E_{1}\) to \(E_{2}\) which is a morphism of representations from \(\varphi_{1}\) into \(\varphi_{2}\) (I, p. 11) and which satisfies \(u(\varphi_{1}(a)x_{1}) = \varphi_{2}(a)x_{2}\) for all \(a \in A\) ;
\item
  The linear map u is isometric;
\item
  If \(\left(\mathrm{E}_{2},x_{2},\varphi_{2}\right)\) is cyclic, then u is an isomorphism;
\item
  If \(\left(\mathrm{E}_2,x_2,\varphi_2\right)\) is cyclic and if A is unital, then \(u(x_{1}) = x_{2}\) .
\end{enumerate}

For \(j \in \{1,2\}\) , define \(\gamma_j \colon A \to E_j\) by \(\gamma_j(a) = \varphi_j(a)x_j\) for all \(a \in A\) . By definition, the pairs \((E_j, \gamma_j)\) are Hilbert realizations of the universally positive kernel \(f\) on \(A\) defined by \((a,b) \mapsto \lambda(a^*b)\) . The Hilbert realization \((E_1, \gamma_1)\) is cyclic; by prop. 1 of V, p. 434, there therefore exists a unique continuous linear map \(u \colon E_1 \to E_2\) such that \(\gamma_2 = u \circ \gamma_1\) , and this map is isometric. Moreover, if \((E_2, x_2, \varphi_2)\) is also cyclic, then \(u\) is an isomorphism. To prove \(a)\) , \(b)\) and \(c)\) , it suffices to prove that \(u\) is a morphism of representations from \(\varphi_1\) into \(\varphi_2\) .

Let \(a \in \mathrm{A}\) . For every \(b \in \mathrm{A}\) , we have

\[
\begin{array}{c} (u \circ \varphi_ {1} (a)) (\gamma_ {1} (b)) = (u \circ \varphi_ {1} (a b)) x _ {1} = (u \circ \gamma_ {1}) (a b) \\ = \gamma_ {2} (a b) = \varphi_ {2} (a) (\gamma_ {2} (b)) = \varphi_ {2} (a) (u (\gamma_ {1} (b))), \end{array}
\]

therefore the continuous linear maps \(u \circ \varphi_1(a)\) and \(\varphi_2(a) \circ u\) coincide on the subspace of \(E_1\) generated by the image of \(\gamma_1\) ; this subspace is dense in \(E_1\) by hypothesis, hence \(u \circ \varphi_1(a) = \varphi_2(a) \circ u\) .

Let us finally prove \(d\) ). Suppose therefore that A is unital and that \((\mathrm{E}_2, x_2, \varphi_2)\) is cyclic. There exists a cyclic Hilbert-space realization \((\mathrm{E}_3, x_3, \varphi_3)\) of \(\lambda\) such that \(\varphi_3\) is a unital morphism (prop. 5). Let \(j \in \{1, 2\}\) . Applying the preceding to \((\mathrm{E}_j, x_j, \varphi_j)\) and \((\mathrm{E}_3, x_3, \varphi_3)\) , one sees that there exists an isometric isomorphism \(\widetilde{u}_j\) from \(\mathrm{E}_j\) into \(\mathrm{E}_3\) such that \(\widetilde{u}_j \circ \varphi_j(a) = \varphi_3(a) \circ \widetilde{u}_j\) for all \(a \in \mathrm{A}\) . Taking \(a = 1_{\mathrm{A}}\) , one sees that \(\varphi_j\) is unital. One then has \(x_2 = \varphi_2(1_{\mathrm{A}})x_2 = u(\varphi_1(1_{\mathrm{A}})x_1) = u(x_1)\) .

Remark. --- Keep the notations of the proposition. It is possible that \(u(x_{1})\) is different from \(x_{2}\) (exercise 3, b) of V, p. 493). However, the triple \((\mathrm{E}_{2}, u(x_{1}), \varphi_{2})\) is also a Hilbertian realization of \(\lambda\) .

\subsection{Representations of stellar algebras}

Let A be a stellar algebra. Denote by \(A_{+}\) the closed convex cone of positive elements of A (def. 6 of I, p. 115). A linear form \(\lambda\) on A is positive if and only if \(\lambda(\mathrm{A}_{+}) \subset \mathbf{R}_{+}\) (I, p. 118, th. 2).

PROPOSITION 7. --- Let A be a stellar algebra. Every positive linear form \(\lambda\) on A is continuous.

Let us first show that \(\lambda\) is bounded on the intersection of \(\mathrm{A}_{+}\) and the unit ball of A. Suppose this is not the case. Then there exists a sequence \((x_{n})_{n\geqslant 1}\) in \(\mathrm{A}_{+}\) such that \(\| x_n\| \leqslant 1\) and \(\lambda (x_n)\geqslant n\) for every integer \(n\geqslant 1\) . The series with general term \(n^{-2}x_{n}\) converges to an element \(x\) of A (cf. TG, IV, p. 33, example 3). For every integer \(\mathrm{N}\geqslant 1\) , since

\[
\sum_ {n = \mathrm{N} + 1} ^ {+ \infty} \frac {1}{n ^ {2}} x _ {n} \in \mathrm{A} _ {+},
\]

(I, p. 116, prop. 14) it follows that

\[
\sum_ {n = 1} ^ {\mathrm{N}} \frac {1}{n} \leqslant \sum_ {n = 1} ^ {\mathrm{N}} \frac {1}{n ^ {2}} \lambda (x _ {n}) = \lambda (x) - \lambda \left(\sum_ {n = \mathrm{N} + 1} ^ {+ \infty} \frac {1}{n ^ {2}} x _ {n}\right) \leqslant \lambda (x),
\]

which is absurd (TG, IV, p. 33, example 4).

Since every element of the unit ball of A is a linear combination with coefficients bounded by 1 of at most four positive elements of A of norm \(\leqslant 1\) (I, p. 96, lemma 2 and formula (4) of I, p. 117), we conclude that \(\lambda\) is continuous.

There exist involutive Banach algebras which admit positive linear forms that are not continuous (exercise 3, a) of V, p. 493).

PROPOSITION 8. --- Let A be a stellar algebra and let a be a nonzero element of A. There exists a positive linear form \(\lambda \in \mathrm{A}_{+}^{\prime}\) such that \(\lambda(a^{*}a) > 0\) .

Consider the real Banach space \(A_{h}\) of Hermitian elements of A. The set \(A_{+}\) is a closed convex pointed salient cone in \(A_{h}\) (prop. 14 of I, p. 116). The element \(a^{*}a\) is positive (th. 2 of I, p. 118) and nonzero, hence the Hermitian element \(-a^{*}a\) is not positive. By EVT, II, p. 42, cor. 5, there exists a continuous real linear form \(\lambda\in A_{h}^{\prime}\) such that \(\lambda(-a^{*}a)<0\) and \(\lambda(\mathrm{A}_{+})\subset\mathbf{R}_{+}\) . The linear form \(\lambda\) extends to a Hermitian C-linear form on A (cf. I, p. 98), which is positive and possesses the required property.

THEOREM 2 (Gelfand--Naimark). --- Let A be a stellar algebra. There exists a Hilbert space E and an isometric morphism of involutive algebras \(\varphi\) from A into \(\mathcal{L}(\mathrm{E})\) . If A is unital, there exists such a unital morphism.

For every \(b \in \mathrm{A} - \{0\}\) , let \(\lambda_b\) be a continuous positive linear form on A such that \(\lambda_b(b^*b) > 0\) (prop. 8). Since A admits an approximate unit (I, p. 121, prop. 18), there exists a Hilbertian realization \((\mathrm{E}_b, x_b, \varphi_b)\) of \(\lambda_b\) (prop. 5). If A is unital, one may suppose that \(\varphi_b\) is unital (loc. cit.). We have \(\| \varphi_b(b) \|^2 = \lambda_b(b^*b) \neq 0\) .

Let E be the external Hilbert sum of the spaces \(E_{b}\) for \(b\) belonging to \(A - \{0\}\) . For every \(a \in A\) , denote by \(\varphi(a) \in \mathcal{L}(E)\) the unique continuous linear map whose restriction to \(E_{b}\) coincides with \(\varphi_{b}(a)\) for all \(b \in A - \{0\}\) . The map \(a \mapsto \varphi(a)\) is a morphism of involutive algebras; it is injective, hence isometric (prop. 9 of I, p. 112), and satisfies \(\varphi(1) = 1_{\mathrm{E}}\) if A is unital. This concludes the proof.

*The category whose objects are the stellar algebras and whose morphisms are the morphisms of involutive algebras is therefore equivalent to the category whose objects are the closed involutive subalgebras of the endomorphism algebras of Hilbert spaces, and whose morphisms are the morphisms of involutive algebras.*

\subsection{Functions of positive type on a topological group}

In this number, G is a topological group whose identity element is denoted by e. The Hilbert spaces considered are complex.

DEFINITION 4. --- A continuous function \(\varphi\in\mathcal{C}(\mathrm{G})\) is said to be of positive type on G if the function \(f\) defined by \(f(g,h)=\varphi(g^{-1}h)\) on \(\mathrm{G}\times\mathrm{G}\) is a universally positive kernel on \(\mathrm{G}\) .

We denote by \(\operatorname{Pos}(\mathbf{G})\) the set of functions of positive type on G and we denote by \(\operatorname{Pos}_{1}(\mathbf{G})\) (resp. \(\operatorname{Pos}_{\leqslant1}(\mathbf{G})\) ) the subset of \(\varphi\in\operatorname{Pos}(\mathbf{G})\) such that \(\varphi(e)=1\) (resp. such that \(\varphi(e)\leqslant1\) ).

Example. --- Let \(\varrho\) be a unitary representation of G in a Hilbert space E and \(x \in \mathrm{E}\) . Let \(\varphi\) be the diagonal matrix coefficient defined by \(\varphi(g) = \langle x \mid \varrho(g)x \rangle\) for all \(g \in \mathrm{G}\) ; the function \(\varphi\) is continuous. It follows that

\[
\varphi (g ^ {- 1} h) = \langle x | \varrho (g) ^ {*} \varrho (h) x \rangle = \langle \varrho (g) x | \varrho (h) x \rangle
\]

for all \((g,h)\in\mathrm{G}\times\mathrm{G}\) . Consequently \(\varphi\in\mathrm{Pos}(\mathrm{G})\) (theorem 1 of V, p. 432, (iv)). We have \(\varphi\in\mathrm{Pos}_{1}(\mathrm{G})\) if and only if \(\|x\|=1\) .

DEFINITION 5. --- Let \(\varphi\) be a function of positive type on G. A Hilbert realization of \(\varphi\) is a pair \((\varrho, x)\) where \(\varrho\) is a unitary representation of G in a Hilbert space E and \(x \in \mathrm{E}\) such that \(\varphi(g) = \langle x \mid \varrho(g)x \rangle\) for all \(g \in \mathrm{G}\) .

If x is a cyclic vector of \(\varrho\) , one says that it is a cyclic Hilbert realization of \(\varphi\) .

PROPOSITION 9. --- Let \(\varphi\in\mathcal{C}(\mathrm{G})\) . The following conditions are equivalent:

\begin{enumerate}
\def\labelenumi{(\roman{enumi})}
\item
  The function \(\varphi\) is of positive type on G;
\item
  There exists a cyclic Hilbert realization of \(\varphi\) ;
\item
  There exists a Hilbert realization of \(\varphi\) .
\end{enumerate}

Condition (ii) implies condition (iii), and condition (iii) implies (i) by the example above.

Let us finally prove that (i) implies (ii). Let (E, \(\gamma\) ) be a cyclic Hilbert realization of the universally positive kernel defined by \(f(g,h)=\varphi(g^{-1}h)\) for \((g,h)\in\mathrm{G}\times\mathrm{G}\) . Let \(k\in\mathrm{G}\) . The continuous function \(\gamma_{k}:g\mapsto\gamma(kg)\) on G satisfies

\[
\langle \gamma_ {k} (g) | \gamma_ {k} (h) \rangle = \langle \gamma (k g) | \gamma (k h) \rangle = f (k g, k h) = \varphi ((k h) ^ {- 1} k g) = f (g, h)
\]

for all \((g,h)\in\mathrm{G}\times\mathrm{G}\) , hence \((\mathrm{E},\gamma_{k})\) is a Hilbert realization of f. By prop. 1 of V, p. 434, there exists a unique unitary operator \(\varrho(k)\) in \(\mathcal{L}(\mathrm{E})\) such that \(\gamma_{k}=\varrho(k)\circ\gamma\) . For every \(g\in G\) and every \((k_{1},k_{2})\in\mathrm{G}\times\mathrm{G}\) , it follows that

\[
\varrho (k _ {1} k _ {2}) (\gamma (g)) = \gamma (k _ {1} k _ {2} g) = \varrho (k _ {1}) (\varrho (k _ {2}) (\gamma (g))),
\]

whence \(\varrho(k_1k_2) = \varrho(k_1)\varrho(k_2)\) since the image of \(\gamma\) is a total subset of E.

Let \(k \in G\) . For every \(g \in G\) , we have \(\varrho(g)(\gamma(k)) = \gamma(gk)\) , and the map from G to E defined by \(g \mapsto \varrho(g)(\gamma(k))\) is therefore continuous. Since the endomorphism \(\varrho(g)\) is unitary, one deduces that \(\varrho\) is a unitary representation of G in E (lemma 4 of V, p. 380). Then set \(x = \gamma(e)\) . The set of vectors \(\varrho(g)x = \varrho(g)(\gamma(e)) = \gamma(g)\) for \(g\) ranging over G is total in E, hence x is a cyclic vector of \(\varrho\) . Since

\[
\langle x \mid \varrho (g) x \rangle = f (e, g) = \varphi (g)
\]

for all \(g \in G\) , the pair \((\varrho, x)\) is a cyclic Hilbert realization of \(\varphi\) .

PROPOSITION 10. — Let \(\varphi\) be a function of positive type on G and let \((\varrho_{1}, x_{1})\) and \((\varrho_{2}, x_{2})\) be Hilbert realizations of \(\varphi\) , the Hilbert representation \((\varrho_{1}, x_{1})\) being cyclic. There exists a unique isometric G-morphism \(u\) from \(\varrho_{1}\) into \(\varrho_{2}\) such that \(u(x_{1}) = x_{2}\) . If \((\varrho_{2}, x_{2})\) is also cyclic, then \(u\) is an isomorphism.

For \(1 \leqslant j \leqslant 2\) , let \(\mathrm{E}_j\) be the space of \(\varrho_j\) and \(\gamma_j\) the function on \(\mathbf{G}\) defined by \(\gamma_j(g) = \varrho_j(g)x_j\) for all \(g \in \mathbf{G}\) . The pairs \((\mathrm{E}_1, \gamma_1)\) and \((\mathrm{E}_2, \gamma_2)\) are Hilbertian realizations of the function of positive type \((g, h) \mapsto \varphi(g^{-1}h)\) , and \((\mathrm{E}_1, \gamma_1)\) is a cyclic Hilbertian realization. By prop. 1 of V, p. 434, there exists a unique isometric linear map \(u: \mathrm{E}_1 \to \mathrm{E}_2\) such that \(\gamma_2 = u \circ \gamma_1\) . In particular, we have \(x_2 = \gamma_2(e) = u(\gamma_1(e)) = u(x_1)\) . Moreover, for all \(g\) and \(h\) in \(\mathbf{G}\) , we have

\[
\varrho_ {2} (g) u \left(\gamma_ {1} (h)\right) = \varrho_ {2} (g) \gamma_ {2} (h) = \gamma_ {2} (g h) = u \left(\gamma_ {1} (g h)\right) = u \left(\varrho_ {1} (g) \gamma_ {1} (h)\right),
\]

and since the set of elements \(\gamma_{1}(h)\) for \(h \in G\) is total in \(E_{1}\) , this means that u is a G-morphism. By loc. cit., it is an isometric isomorphism if \((\varrho_{2}, x_{2})\) is also cyclic.

PROPOSITION 11. — Let \(\varphi \in \mathrm{Pos}_{1}(\mathrm{G})\) and \((\varrho, x)\) be a cyclic Hilbertian realization of \(\varphi\) . Then \(\varrho\) is an irreducible representation of G if and only if \(\varphi\) is an extreme point (EVT, II, p. 57, def. 1) of \(\mathrm{Pos}_{1}(\mathrm{G})\) .

Let E be the space of \(\varrho\) . Suppose first that \(\varrho\) is not irreducible. Let F be a closed subspace of E stable under \(\varrho\) , nonzero and different from E. Let us write \(x = x_1 + x_2\) , where \(x_1 \in \mathrm{F}\) and \(x_2 \in \mathrm{F}^\circ\) . We then have \(1 = \|x\|^2 = \|x_1\|^2 + \|x_2\|^2\) . The subrepresentation of E generated by \(x_1\) is contained in F, hence \(x_1 \neq x\) since \(x\) is a cyclic vector of \(\varrho\) , whence \(x_2 \neq 0\) . Similarly, one verifies that \(x_1 \neq 0\) .

For \(j = 1\) and \(j = 2\) , let \(\varphi_j\) be the continuous function on G such that

\[
\varphi_ {j} (g) = \frac {1}{\| x _ {j} \| ^ {2}} \langle x _ {j} | \varrho (g) x _ {j} \rangle
\]

for all \(g \in G\) . We have \(\varphi_{j} \in \mathrm{Pos}_{1}(G)\) (V, p. 443, example). Since \(\varphi = \|x_{1}\|^{2}\varphi_{1} + \|x_{2}\|^{2}\varphi_{2}\) , it suffices to check that \(\varphi_{1} \neq \varphi_{2}\) to show that \(\varphi\) is not an extreme point of \(\mathrm{Pos}_{1}(G)\) ; it suffices for that to prove that \(\varphi \neq \varphi_{1}\) .

Let us reason by contradiction and suppose that \(\varphi = \varphi_1\) . Since \(\langle x_1 | \varrho(g)x_2 \rangle = 0\) for all \(g \in G\) , we would have

\[
\frac {1}{\| x _ {1} \| ^ {2}} \langle x _ {1} | \varrho (g) x \rangle = \frac {1}{\| x _ {1} \| ^ {2}} \langle x _ {1} | \varrho (g) x _ {1} \rangle = \varphi_ {1} (g) = \varphi (g) = \langle x | \varrho (g) x \rangle
\]

for all \(g \in G\) , whence \(\langle x_{1} | y \rangle = \langle \|x_{1}\|^{2}x | y \rangle\) for every element y of the vector subspace of E generated by the elements \(\varrho(g)x\) for \(g \in G\) , therefore, for all \(y \in E\) , since x is a cyclic vector of \(\varrho\) . This would imply that \(x_{1} = \|x_{1}\|^{2}x\) is also a cyclic vector of \(\varrho\) , which is a contradiction, whence the assertion.

Suppose now that \(\varrho\) is irreducible and let us prove that \(\varphi\) is an extreme point of \(\mathrm{Pos}_1(\mathrm{G})\) . Let \(\varphi_1 \neq \varphi_2\) be elements of \(\mathrm{Pos}_1(\mathrm{G})\) and \(t_1, t_2 \in [0,1]\) such that \(t_1 + t_2 = 1\) and \(\varphi = t_1\varphi_1 + t_2\varphi_2\) . For \(j \in \{1,2\}\) , let \((\varrho_j, x_j)\) be a cyclic Hilbertian realization of \(\varphi_j\) , and let \(\mathrm{E}_j\) be the space of \(\varrho_j\) . Let \(x_3 = t_1^{1/2}x_1 + t_2^{1/2}x_2\) . Then \((\varrho_1 \oplus \varrho_2, x_3)\) is a Hilbert realization of \(\varphi\) . Since \((\varrho, x)\) is cyclic, there exists an isometric G-morphism \(u: \mathrm{E} \to \mathrm{E}_1 \oplus \mathrm{E}_2\) such that \(u(x) = x_3\) (prop. 10).

Let \(j = 1\) or \(j = 2\) . Since \(\varrho\) is irreducible, there exists \(\lambda_j \geqslant 0\) such that the G-morphism \(u_j = \mathrm{pr}_j \circ u\) from E into \(E_j\) satisfies

\[
\langle u _ {j} (y) \mid u _ {j} (y ^ {\prime}) \rangle = \lambda_ {j} \langle y \mid y ^ {\prime} \rangle
\]

for all \(y\) and \(y'\) in E (cor. 5 of V, p. 388, if \(u_j \neq 0\) , and one may take \(\lambda_j = 0\) otherwise). For every \(g \in G\) , it follows that

\[
\begin{array}{r} t _ {j} \varphi_ {j} (g) = \langle t _ {j} ^ {1 / 2} x _ {j} | \varrho_ {j} (g) (t _ {j} ^ {1 / 2} x _ {j}) \rangle = \langle u _ {j} (x) | \varrho_ {j} (g) (u _ {j} (x)) \rangle \\ = \langle u _ {j} (x) | u _ {j} (\varrho (g) x) \rangle = \lambda_ {j} \varphi (g). \end{array}
\]

Since \(\varphi_{j}(e) = \varphi(e) = 1\) , we deduce that \(\varphi_{j} = \varphi\) if \(t_{j} \neq 0\) . The hypothesis \(\varphi_{1} \neq \varphi_{2}\) therefore implies that \(t_{1}\) or \(t_{2}\) must be zero, which proves that \(\varphi\) is an extreme point of \(\mathrm{Pos}_{1}(\mathrm{G})\) .

Lemma 4. --- Let \(\varphi\in\mathrm{Pos}(\mathbf{G})\) . The function \(\varphi\) is bounded on \(\mathbf{G}\) and we have \(\|\varphi\|_{\infty}=\varphi(e)\) . Moreover, we have

\[
| \varphi (g ^ {- 1} h) - \varphi (h) | \leqslant \sqrt {2 \varphi (e) (\varphi (e) - \mathcal {R} (\varphi (g)))}\tag{1}
\]

for all \((g,h)\in \mathbf{G}\times \mathbf{G}\) .

Let \((\varrho, x)\) be a Hilbert realization of \(\varphi\) . We have \(\varphi(e) = \|x\|^2\) . For every \(g \in G\) , it follows therefore that \(|\varphi(g)| = |\langle x | \varrho(g)x\rangle| \leqslant \varphi(e)\) by the Cauchy-Schwarz inequality. This proves the first assertion.

For every \((g,h)\in \mathrm{G}\times \mathrm{G}\) , we have

\[
\varphi (g ^ {- 1} h) - \varphi (h) = \langle x | \varrho (g ^ {- 1} h) x \rangle - \langle x | \varrho (h) x \rangle = \langle \varrho (g) x - x | \varrho (h) x \rangle
\]

since \(\varrho\) is unitary, whence

\[
| \varphi (g ^ {- 1} h) - \varphi (h) | \leqslant \| \varrho (g) x - x \| \| \varrho (h) x \| \leqslant \varphi (e) ^ {1 / 2} \| \varrho (g) x - x \|.
\]

One concludes by observing that

\[
\| \varrho (g) x - x \| ^ {2} = 2 \| x \| ^ {2} - 2 \mathcal {R} (\langle x | \varrho (g) x \rangle) = 2 (\varphi (e) - \mathcal {R} (\varphi (g))).
\]

Remark. --- The sets \(\operatorname{Pos}(G)\) (resp. \(\operatorname{Pos}_{1}(G)\) and \(\operatorname{Pos}_{\leqslant1}(G)\) ) are convex self-adjoint subsets of the stellar algebra \(\mathcal{C}_{b}(G)\) . They are closed in the space \(\mathcal{C}_{b}(G)\) equipped with the topology of simple convergence. The set \(\operatorname{Pos}(G)\) is a convex cone with vertex 0 in the real Banach space \(\mathcal{C}_{b}(G)\) .

\subsection{Unitary dual of a locally compact group}

Let G be a locally compact topological group. One equips G with a left Haar measure \(\mu\) . For \(p \in [1, +\infty]\) , we denote by \(\mathcal{L}^{p}(\mathrm{G})\) (resp. \(\mathrm{L}^{p}(\mathrm{G})\) ) the space \(\mathcal{L}_{\mathbf{C}}^{p}(\mathrm{G}, \mu)\) (resp. the space \(\mathrm{L}_{\mathbf{C}}^{p}(\mathrm{G}, \mu)\) ). We identify the space \(\mathcal{C}_{b}(\mathrm{G})\) with its image in \(\mathrm{L}^{\infty}(\mathrm{G})\) .

Denote by \(\Delta\) the modular function of G. Recall that \(L^1(G)\) is an involutive Banach algebra whose involution is induced, by passage to quotients, by the map \(f \mapsto f^*\) where \(f^*(y) = \Delta^{-1}(y)\overline{f(y^{-1})}\) for all \(f \in \mathcal{L}^1(G)\) and \(y \in G\) (cf. I, p. 99, example 4). The Banach algebra \(L^1(G)\) admits an approximate unit according to INT, VIII, p. 172, §4, no. 7, prop. 20.

Let \(\varrho\) be a unitary representation of G in a complex Hilbert space E. The map \(f\mapsto\varrho(f)\) is a continuous morphism of involutive algebras from L \(^{1}\) (G) to \(\mathcal{L}\) (E) (lemma 1 of V, p. 401); it is a non-degenerate representation of L \(^{1}\) (G) in E (INT, VIII, p. 139, § 2, no. 7, prop. 10, (i)).

PROPOSITION 12. --- Let E be a complex Hilbert space and \(\widetilde{\pi}\) a morphism of involutive algebras from L \(^{1}\) (G) to \(\mathcal{L}\) (E) . If the representation \(\widetilde{\pi}\) is nondegenerate, then there exists a unique unitary representation \(\pi\) of G in E such that \(\widetilde{\pi}(f)=\pi(f)\) for all \(f\in\mathrm{L}^{1}(\mathrm{G})\) .

Uniqueness follows from INT, VIII, p. 139, § 2, n° 7, cor. 3 of lemma 4.

Let us prove the existence of \(\pi\) . We have \(\| \widetilde{\pi}\| \leqslant 1\) by prop. 2 of I, p. 104. Let F be the subspace of E generated by the vectors of the form \(\widetilde{\pi}(f)x\) for \(f\) in \(\mathrm{L}^1 (\mathrm{G})\) and \(x\) in E; it is dense in E since the representation \(\widetilde{\pi}\) is nondegenerate.

For every compact neighborhood V of e, let \(\varphi_{V}\) be a continuous positive function with support contained in V such that \(\int\varphi_{V}\mu=1\) . Let B be a base of the filter of compact neighborhoods of e.

Let \(g\) be in \(G\) and \(f\) be in \(L^1(G)\) . We have \((\varphi_V * \varepsilon_g) * f \to \varepsilon_g * f\) in \(L^1(G)\) along the filter of sections of \(\mathfrak{B}\) (INT, VIII, p. 172, § 4, no. 7, prop. 20), hence \(\widetilde{\pi}(\varphi_V * \varepsilon_g)\widetilde{\pi}(f)\) converges to \(\widetilde{\pi}(\varepsilon_g * f)\) in \(\mathcal{L}(E)\) . This implies that \(\widetilde{\pi}(\varphi_V * \varepsilon_g)\) converges simply along the filter of sections of \(\mathfrak{B}\) to a linear map \(\pi(g)\) from \(F\) to \(F\) ; this map is continuous by EVT, III, p. 26, cor. 3, since \(\| \widetilde{\pi}(\varphi_V * \varepsilon_g) \| \leqslant 1\) for every compact neighborhood \(V\) of \(e\) . There therefore exists a unique endomorphism of \(E\) which induces \(\pi(g)\) by passage to the subspaces. We still denote by \(\pi(g)\) this endomorphism; we have \(\| \pi(g) \| \leqslant 1\) .

Let \(f \in \mathrm{L}^1(\mathrm{G})\) . Since \(\widetilde{\pi}(\varphi_{\mathrm{V}} * \varepsilon_g) \widetilde{\pi}(f)\) converges to \(\widetilde{\pi}(\varepsilon_g * f)\) , we have \(\pi(g) \widetilde{\pi}(f) = \widetilde{\pi}(\varepsilon_g * f)\) .

For \(g = e\) , this relation shows that \(\pi(e)\) is the identity on \(F\) , hence \(\pi(e) = 1_{E}\) . Let \(g_1\) and \(g_2\) be in \(G\) . We have

\[
\pi (g _ {1}) \pi (g _ {2}) \widetilde {\pi} (f) = \pi (g _ {1}) \widetilde {\pi} (\varepsilon_ {g _ {2}} * f) = \widetilde {\pi} (\varepsilon_ {g _ {1} g _ {2}} * f) = \pi (g _ {1} g _ {2}) \widetilde {\pi} (f),
\]

whence \(\pi(g_1)\pi(g_2) = \pi(g_1g_2)\) on F, and hence on E. This proves that the map \(g \mapsto \pi(g)\) is a representation of G in E. Since \(\|\pi(g)\| \leqslant 1\) and \(\|\pi(g)^{-1}\| = \|\pi(g^{-1})\| \leqslant 1\) , the endomorphisms \(\pi(g)\) of E are isometric, hence unitary (EVT, V, p. 40, prop. 3).

Let \(f \in \mathrm{L}^{1}(\mathrm{G})\) and \(x \in E\) . The map \(g \mapsto \varepsilon_{g} * f\) from G into \(\mathrm{L}^{1}(\mathrm{G})\) is continuous (INT, VIII, p. 136, § 2, no. 5 and p. 144, § 3, no. 2, formula (5)), and \(\widetilde{\pi}\) is continuous, hence the map \(g \mapsto \widetilde{\pi}(\varepsilon_{g} * f)x = \pi(g)\widetilde{\pi}(f)x\) is continuous at \(g = e\) . Since F is dense in E, the representation \(\pi\) is unitary (lemma 4 of V, p. 380).

Let \(f_1\) and \(f_2\) be in \(L^1(G)\) . According to INT, VIII, p. 127, §1, no. 5, prop. 7, we have the relation

\[
f _ {1} * f _ {2} = \int_ {\mathrm{G}} f _ {1} (g) (\varepsilon_ {g} * f _ {2}) d \mu (g)
\]

in L \(^{1}\) (G), whence

\[
\begin{array}{l} \widetilde {\pi} (f _ {1}) \widetilde {\pi} (f _ {2}) = \widetilde {\pi} (f _ {1} * f _ {2}) = \int_ {\mathrm{G}} f _ {1} (g) \widetilde {\pi} (\varepsilon_ {g} * f _ {2}) d \mu (g) \\ \qquad = \int_ {\mathrm{G}} f _ {1} (g) \pi (g) \widetilde {\pi} (f _ {2}) d \mu (g) = \Big (\int_ {\mathrm{G}} f _ {1} (g) \pi (g) d \mu (g) \Big) \widetilde {\pi} (f _ {2}), \end{array}
\]

using INT, VI, p. 9, § 1, no. 1, prop. 1. It follows that

\[
\widetilde {\pi} (f) = \int_ {\mathrm{G}} f (g) \pi (g) d \mu (g) = \pi (f)
\]

for all \(f \in \mathrm{L}^{1}(\mathrm{G})\) . The proposition is proved.

PROPOSITION 13. --- Let \(\varphi\in\mathrm{L}^{\infty}(\mathrm{G})\) . Denote by \(\lambda_{\varphi}\) the continuous linear form \(f\mapsto\langle f,\varphi\rangle\) on \(\mathrm{L}^{1}(\mathrm{G})\) . Then \(\varphi\) is the class of a function of positive type on G if and only if \(\lambda_{\varphi}\) is a continuous positive linear form on \(\mathrm{L}^{1}(\mathrm{G})\) .

Suppose that \(\lambda_{\varphi}\) is a continuous positive linear form on \(L^{1}(G)\) . Let \((E, x, \widetilde{\pi})\) be a cyclic Hilbertian realization of \(\lambda\) (prop. 5 of V, p. 438).

Let \(\pi\) be a unitary representation of G in E such that \(\pi(f) = \widetilde{\pi}(f)\) for all \(f \in \mathrm{L}^1(\mathrm{G})\) (prop. 12). We then have \(\lambda_{\varphi}(f) = \langle x | \pi(f)x \rangle\) for all \(f \in \mathrm{L}^1(\mathrm{G})\) .

Since

\[
\pi (f) = \int_ {\mathrm{G}} f (g) \pi (g) d \mu (g)
\]

for all \(f\in \mathrm{L}^1 (\mathrm{G})\) , it follows that

\[
\int_ {\mathrm{G}} f (g) \varphi (g) d \mu (g) = \lambda_ {\varphi} (f) = \langle x | \pi (f) x \rangle = \int_ {\mathrm{G}} f (g) \langle x | \pi (g) x \rangle d \mu (g)
\]

for all \(f \in \mathrm{L}^{1}(\mathrm{G})\) . Thus, \(\varphi\) is the class in \(\mathrm{L}^{\infty}(\mathrm{G})\) of the function on G defined by \(g \mapsto \langle x | \pi(g)x \rangle\) , which is a function of positive type on G (Prop. 9 of V, p. 443).

Conversely, let \(\varphi\in\mathrm{Pos}(\mathrm{G})\) . Let \(f\in\mathcal{K}(\mathrm{G})\) . We then have

\[
\begin{array}{c} \lambda_ {\varphi} (f ^ {*} * f) = \int_ {G} \varphi (y) \Big (\int_ {G} \Delta (z) ^ {- 1} \overline {{f (z ^ {- 1})}} f (z ^ {- 1} y) d \mu (z) \Big) d \mu (y) \\ = \int_ {G} \varphi (y) \Big (\int_ {G} \overline {{f (z)}} f (z y) d \mu (z) \Big) d \mu (y) \end{array}
\]

The continuous function on G × G defined by \((z, y) \mapsto \varphi(y)\overline{f(z)}f(zy)\) is bounded and

\[
\begin{array}{c} \int_ {\mathrm{G}} ^ {*} | \varphi (y) | \Big (\int_ {\mathrm{G}} ^ {*} | \overline {{f (z)}} f (z y) | d \mu (z) \Big) d \mu (y) \\ \leqslant \varphi (e) \int_ {\mathrm{G}} ^ {*} \int_ {\mathrm{G}} ^ {*} | f (z) | | f (z y) | d \mu (z) d \mu (y) \\ = \varphi (e) \Big (\int_ {\mathrm{G}} ^ {*} | f (z) | d \mu (z) \Big) ^ {2} \end{array}
\]

By INT, V, p. 94, § 8, no. 3, prop. 8, we therefore deduce from the Lebesgue–Fubini theorem (INT, V, p. 96, § 8, no. 4, th. 1) that

\[
\begin{array}{c} \lambda_ {\varphi} (f ^ {*} * f) = \int_ {\mathrm{G}} \overline {{f (z)}} \Bigl (\int_ {\mathrm{G}} \varphi (y) f (z y) d \mu (y) \Bigr) d \mu (z) \\ = \int_ {\mathrm{G} \times \mathrm{G}} \overline {{f (z)}} f (y) \varphi (z ^ {- 1} y) d (\mu \otimes \mu) (z, y). \end{array}
\]

This implies that \(\lambda_{\varphi}(f^{*} * f) \geqslant 0\) by Theorem 1, (i) of V, p. 432 applied to the measure induced by \(\mu\) on the support of \(f\) (INT, IV, p. 186, § 5, n° 7, def. 4). From this one deduces by continuity that \(\lambda_{\varphi}(f^{*} * f) \geqslant 0\) for all \(f \in \mathrm{L}^1(\mathrm{G})\) .

The spaces \(\mathrm{L}^{\infty}(\mathrm{G})\) and \(\mathcal{C}_{b}(\mathrm{G})\) are equipped with their Banach space topologies, and we denote the norm by \(f \mapsto \|f\|_{\infty}\) . We shall call here the weak topology on \(\mathrm{L}^{\infty}(\mathrm{G})\) , \(\mathcal{C}_{b}(\mathrm{G})\) or \(\operatorname{Pos}(\mathrm{G})\) the topology induced by the weak topology \(\sigma(\mathrm{L}^{\infty}(\mathrm{G}), \mathrm{L}^{1}(\mathrm{G}))\) .

Since \(L^{\infty}(G)\) is identified with the dual of \(L^{1}(G)\) (INT, V, p. 61, §5, no. 8, th. 4), every closed ball of \(L^{\infty}(G)\) is compact for the weak topology (EVT, III, p. 17, cor. 3).

COROLLARY. --- In \(\mathrm{L}^{\infty}(\mathrm{G})\) equipped with the weak topology, the set \(\operatorname{Pos}(\mathrm{G})\) is closed, and the set \(\operatorname{Pos}_{\leqslant 1}(\mathrm{G})\) is compact.

The first assertion follows from the proposition, and the second then follows from EVT, III, loc. cit.

In general, \(\mathrm{Pos}_{1}(\mathrm{G})\) is not compact for the weak topology, as shown by the example of the group R (exercise 10 of V, p. 497).

PROPOSITION 14. --- The set of extreme points of \(\mathrm{Pos}_{\leqslant 1}(\mathrm{G})\) is equal to the union of the set of extreme points of \(\mathrm{Pos}_1(\mathrm{G})\) and of the zero function. Moreover, the closed convex hull of the set of extreme points of \(\mathrm{Pos}_1(\mathrm{G})\) contains \(\mathrm{Pos}_1(\mathrm{G})\) .

By the corollary of Proposition 13, the set \(\mathrm{Pos}_{\leqslant 1}(\mathrm{G})\) is a cap of the pointed convex cone \(\mathrm{Pos}(\mathrm{G})\) , which is defined by the gauge \(\varphi \mapsto \varphi(e)\) (EVT, II, p. 61, def. 3 and prop. 4). Its extreme points are therefore the zero function and the elements \(\varphi\) of \(\mathrm{Pos}_1(\mathrm{G})\) belonging to the extreme generators of \(\mathrm{Pos}(\mathrm{G})\) (EVT, II, p. 62, cor. 1). These are the extreme points of \(\mathrm{Pos}_1(\mathrm{G})\) (EVT, II, p. 61, prop. 3). The first assertion follows from this.

Let us prove the second assertion. Let \(\varphi\in\mathrm{Pos}_{1}(\mathrm{G})\) . Since the set \(\mathrm{Pos}_{\leqslant 1}(\mathrm{G})\) is compact for the weak topology (corollary of Prop. 13), there exists a filter \(\mathfrak{F}\) on the convex hull of the extreme points of \(\mathrm{Pos}_{\leqslant 1}(\mathrm{G})\) which converges to \(\varphi\) (EVT, II, p. 59, th. 1). As the closed balls of the Banach space \(\mathrm{L}^{\infty}(\mathrm{G})\) are closed for the weak topology and since \(\|\varphi\|_{\infty}=1\) , we have \(\lim_{\psi,\mathfrak{F}}\|\psi\|_{\infty}=1\) (indeed, otherwise there would exist a real number \(c<1\) and a filter \(\mathfrak{G}\) on the closed ball of radius \(c\) in \(\mathrm{L}^{\infty}(\mathrm{G})\) which would be finer than \(\mathfrak{F}\) , which would imply that \(\varphi\) belongs to this closed ball).

According to the description of the extreme points of \(\mathrm{Pos}_{\leqslant1}(\mathrm{G})\) , every element \(\psi\) of the convex hull of the extreme points of \(\mathrm{Pos}_{\leqslant1}(\mathrm{G})\) is a function of positive type on G of the form

\[
\psi = \sum_ {i \in \mathrm{I}} t _ {i} \psi_ {i},
\]

where I is a finite set, \(\psi_{i}\) is an extreme point of \(\mathrm{Pos}_{1}(\mathrm{G})\) for all \(i \in I\) , and \(t_{i} \in [0,1]\) . We have \(\sum_{i} t_{i} = \psi(e) = \|\psi\|_{\infty} \leqslant 1\) (lemma 4 of V, p. 446). If \(\psi \neq 0\) , the function

\[
\frac {\psi}{\| \psi \| _ {\infty}} = \sum_ {i \in \mathrm{I}} \frac {t _ {i}}{\psi (e)} \psi_ {i}
\]

therefore belongs to the convex hull of the extreme points of \(\mathrm{Pos}_{1}(\mathrm{G})\) . Since \(\lim_{\psi,\mathfrak{F}}\|\psi\|_{\infty}^{-1}\psi=\varphi\) , we conclude that \(\varphi\) belongs to the closed convex hull of the extreme points of \(\mathrm{Pos}_{1}(\mathrm{G})\) . The assertion follows.

\subsection{Existence of irreducible representations}

We keep the notation from the previous number.

THEOREM 3 (Raikov). --- The weak topology on Pos1(G) coincides with the topology of compact convergence.

We will first prove a few lemmas.

Lemma 5. --- Let X be a locally compact topological space and let \(\nu\) be a positive measure on X. On every bounded subset of \(\mathcal{C}_{b}(X)\) the topology induced by the weak topology \(\sigma(L^{\infty}(X), L^{1}(X))\) is coarser than the topology of compact convergence.

Let B be a bounded subset of \(\mathcal{C}_{b}(\mathrm{X})\) and let \(M \in R_{+}\) such that \(\|\varphi\|_{\infty} \leqslant M\) for all \(\varphi \in B\) . Let \(\psi \in L^{1}(\mathrm{X}, \nu)\) and \(\varepsilon > 0\) . Fix a compact subset K of X such that

\[
\int_ {\mathrm{X-K}} | \psi | d \nu <   \varepsilon .
\]

For all \(\varphi_{1}\) and \(\varphi_{2}\) in B, we then have

\[
\begin{array}{r l} & {\langle \varphi_ {1} - \varphi_ {2}, \psi \rangle = \int_ {\mathrm{X}} \psi (\varphi_ {1} - \varphi_ {2}) d \nu} \\ & {\qquad = \int_ {\mathrm{K}} \psi (\varphi_ {1} - \varphi_ {2}) d \nu + \int_ {\mathrm{X-K}} \psi (\varphi_ {1} - \varphi_ {2}) d \nu ,} \end{array}
\]

whence

\[
| \langle \varphi_ {1} - \varphi_ {2}, \psi \rangle | \leqslant \sup _ {x \in \mathrm{K}} | \varphi_ {1} (x) - \varphi_ {2} (x) | \| \psi \| _ {1} + 2 \mathrm{M} \varepsilon ,
\]

and the lemma follows.

Lemma 6. --- Let \(\psi \in \mathrm{L}^{1}(\mathrm{G})\) . Let B be a bounded subset of the Banach space \(\mathrm{L}^{\infty}(\mathrm{G})\) . The map \(\varphi \mapsto \psi * \varphi\) from B to \(\mathcal{C}_b(\mathrm{G})\) is continuous when B is equipped with the weak topology and \(\mathcal{C}_b(\mathrm{G})\) with the topology of compact convergence.

Let \(\varphi\in\mathrm{L}^{\infty}(\mathrm{G})\) and \(\eta\in\mathrm{L}^{1}(\mathrm{G})\) . The function \(\Delta^{-1}\check{\eta}\) belongs to \(\mathrm{L}^{1}(\mathrm{G})\) and \(\langle\eta,\check{\psi}\rangle=\langle\Delta^{-1}\check{\eta},\psi\rangle\) (cf. INT, VII, p. 19, § 1, n°3, formula (22)). The map \(\varphi\mapsto\check{\varphi}\) is therefore an automorphism of the space \(\mathrm{L}^{\infty}(\mathrm{G})\) endowed with the weak topology.

Let \(\varphi\in\mathrm{L}^{\infty}(\mathrm{G})\) . By INT, VIII, p. 167, § 4, no. 5, prop. 14, the function \(\psi*\varphi\) belongs to \(\mathcal{C}_{b}(\mathrm{G})\) and satisfies

\[
\begin{array}{r l} & {(\psi * \varphi) (g) = \int_ {\mathrm{G}} \psi (h) \varphi (h ^ {- 1} g) d \mu (h)} \\ & {\qquad = \int_ {\mathrm{G}} \psi (g y) \check {\varphi} (y) d \mu (y) = \langle \check {\varphi}, \pmb {\gamma} _ {\mathrm{G}} ^ {(1)} (g ^ {- 1}) \psi \rangle} \end{array}
\]

for all \(g \in G\) . It follows that the linear map \(u: \varphi \mapsto \psi * \varphi\) is a continuous map from the space \(L^{\infty}(G)\) equipped with the weak topology to \(\mathcal{C}_{b}(G)\) endowed with the topology of simple convergence.

Let \(\varphi\in\mathrm{B}\) and \((g,h)\in\mathrm{G}\times\mathrm{G}\) . By the formula above, we have

\[
| u (\varphi) (g) - u (\varphi) (h) | \leqslant \| \check {\varphi} \| _ {\infty} \| (\boldsymbol {\gamma} _ {\mathrm{G}} ^ {(1)} (g ^ {- 1}) - \boldsymbol {\gamma} _ {\mathrm{G}} ^ {(1)} (h ^ {- 1})) \psi \| _ {1}.
\]

As B is bounded and the left regular representation of G in \(L^{1}(G)\) is continuous (no. 4 of V, p. 405), this implies that \(u(B)\) is an equicontinuous subset of \(\mathcal{C}_{b}(G)\) . The assertion then follows from what precedes and from TG, X, p. 16, th. 1.

Lemma 7. --- Let \(\psi \in \mathrm{L}^1 (\mathbf{G})\) such that \(\psi \geqslant 0\) and \(\int \psi = 1\) . Let p denote the semi-norm on \(\mathcal{C}_b(\mathbf{G})\) defined by \(p(\varphi) = |\langle \varphi ,\psi \rangle |\) for all \(\varphi \in \mathcal{C}_b(\mathbf{G})\) . For every \(\varphi \in \mathrm{Pos}_1(\mathbf{G})\) , we have

\[
\| \psi * \varphi - \varphi \| _ {\infty} \leqslant \sqrt {2 p (1 - \varphi)}.
\]

Let \(\varphi\in\mathrm{Pos}_{1}(\mathrm{G})\) and \(x\in\mathrm{G}\) . According to INT, VIII, p. 167, § 4, no. 5, prop. 14, we obtain

\[
\begin{array}{r l} | \psi * \varphi (x) - \varphi (x) | = & \left| \int_ {\mathrm{G}} (\varphi (y ^ {- 1} x) - \varphi (x)) \psi (y) d \mu (y) \right| \\ & \leqslant \int_ {\mathrm{G}} | \varphi (y ^ {- 1} x) - \varphi (x) | \psi (y) d \mu (y) \\ & \leqslant \int_ {\mathrm{G}} \sqrt {2 (1 - \mathcal {R}   \varphi (y))} \psi (y) d \mu (y) \end{array}
\]

by applying the bound (1) from V, p. 446. The Cauchy-Schwarz inequality then implies

\[
\begin{array}{c} \| \psi * \varphi - \varphi \| _ {\infty} \leqslant \sqrt {2} \Bigl (\int_ {\mathrm{G}} (1 - \mathcal {R} (\varphi)) \psi   d \mu \Bigr) ^ {1 / 2} \Bigl (\int_ {\mathrm{G}} \psi   d \mu \Bigr) ^ {1 / 2} \\ \leqslant \sqrt {2} \sqrt {p (1 - \varphi)}, \end{array}
\]

whence the result.

Lemma 8. --- Let K be a topological field, and let E and F be topological vector spaces over K. Let f be a map from E into F and X a subset of E. Let \(x \in X\) . The map f \textbar{} X is continuous at x if, for every neighborhood W of 0 in F, there exist a neighborhood U of x in E and a continuous map g from X to F such that \((f - g)(\mathrm{U} \cap \mathrm{X}) \subset \mathrm{W}\) .

Let \(W_{0}\) be a neighborhood of 0 in F and let \(V_{0}\) be a symmetric neighborhood of 0 in F such that \(V_{0} + V_{0} + V_{0} \subset W_{0}\) . Let \(U_{0}\) be a neighborhood of x in E and let g be a continuous map from X into F such that \((f - g)(\mathrm{U}_{0} \cap \mathrm{X}) \subset \mathrm{V}_{0}\) .

There exists a neighborhood \(U_{1}\) of x in E such that \(g(\mathrm{U}_{1} \cap \mathrm{X}) \subset g(x) + \mathrm{V}_{0}\) . For every \(y \in U_{0} \cap U_{1} \cap X\) , we have

\[
\begin{array}{c} f (y) - f (x) = (f (y) - g (y)) + (g (y) - g (x)) + (g (x) - f (x)) \in V_{0} + V_{0} + V_{0} \subset W_{0}, \end{array}
\]

hence \(f(y) \in f(x) + W_{0}\) , whence the result.

Let us now prove Theorem 3. Denote by \(\mathrm{Pos}_{1}(\mathrm{G})_{f}\) (resp. \(\mathrm{Pos}_{1}(\mathrm{G})_{c}\) ) the set \(\mathrm{Pos}_{1}(\mathrm{G})\) endowed with the weak topology (resp. the topology of compact convergence); likewise, denote by \(\mathcal{C}_{b}(\mathrm{G})_{f}\) (resp. \(\mathcal{C}_{b}(\mathrm{G})_{c}\) ) the space \(\mathcal{C}_{b}(\mathrm{G})\) endowed with the weak topology (resp. the topology of compact convergence).

The identity map from \(\mathrm{Pos}_{1}(\mathrm{G})_{c}\) to \(\mathrm{Pos}_{1}(\mathrm{G})_{f}\) is continuous (Lemma 5). Conversely, denote by \(\iota\) the inclusion of \(\mathrm{Pos}_{1}(\mathrm{G})_{f}\) into \(\mathcal{C}_{b}(\mathrm{G})_{c}\) . Let us prove that \(\iota\) is continuous to conclude the proof.

Let \(\varphi_{1} \in \mathrm{Pos}_{1}(\mathbf{G})\) . To verify that \(\iota\) is continuous at \(\varphi_{1}\) , we apply Lemma 8. Let W be a neighborhood of 0 in \(\mathcal{C}_{b}(\mathbf{G})_{c}\) . It suffices to find a linear map \(u\) from \(\mathcal{C}_{b}(\mathbf{G})\) to \(\mathcal{C}_{b}(\mathbf{G})\) which induces, by passing to subspaces, a continuous map from \(\mathrm{Pos}_{1}(\mathbf{G})_{f}\) to \(\mathcal{C}_{b}(\mathbf{G})_{c}\) , as well as a neighborhood U of 0 in \(\mathrm{Pos}_{1}(\mathbf{G})_{f}\) so that one has \((\iota - u)(\mathrm{U} \cap \mathrm{Pos}_{1}(\mathbf{G})) \subset \mathrm{W}\) .

There exists \(\varepsilon > 0\) such that W contains the set of \(\varphi \in \mathcal{C}_b(\mathrm{G})\) which satisfy \(\| \varphi \|_{\infty} \leqslant \varepsilon\) . Since \(\varphi_1(e) = 1\) , there exists a compact neighborhood V of \(e\) in G such that

\[
\sup _ {x \in \mathrm{V}} | 1 - \varphi_ {1} (x) | \leqslant \frac {\varepsilon^ {2}}{4}.
\]

Let \(\varphi_{\mathrm{V}}\) be the characteristic function of V, and set \(\psi_{\mathrm{V}} = \mu (\mathrm{V})^{-1}\varphi_{\mathrm{V}}\) . The linear map \(u\colon \varphi \mapsto \psi_{\mathrm{V}}*\varphi\) from \(\mathrm{Pos}_1(\mathrm{G})_f\) to \(\mathcal{C}_b(\mathrm{G})_c\) is continuous (Lemma 6).

Denote by \(q_{V}\) the seminorm \(\varphi\mapsto|\langle\varphi,\psi_{V}\rangle|\) on \(\mathcal{C}_{b}(G)\) ; it is continuous for the weak topology. Since \(\psi_{V}\) is zero outside V, it follows that \(q_{\mathrm{V}}(1-\varphi_{1})\leqslant\varepsilon^{2}/4\) ; there therefore exists a neighborhood U of \(\varphi_{1}\) in \(\mathcal{C}_{b}(\mathrm{G})_{f}\) such that \(q_{\mathrm{V}}(1-\varphi)\leqslant\varepsilon^{2}/2\) for all \(\varphi\in U\) .

Let \(\varphi\in\mathrm{U}\cap\mathrm{Pos}_{1}(\mathrm{G})\) . By Lemma 7, we have

\[
\| (\iota - u) (\varphi) \| _ {\infty} \leqslant \sqrt {2 q _ {\mathrm{V}} (1 - \varphi)} \leqslant \varepsilon
\]

hence \((\iota - u)(\mathrm{U}) \subset \mathrm{W}\) . The theorem is proved.

In general, the weak topology on \(\operatorname{Pos}_{\leqslant1}(\mathbf{G})\) does not coincide with the topology of compact convergence, as shown by the example of the group R.

We recall that we denote by \(\widehat{G}\) the set of classes of irreducible unitary representations of \(G\) (V, p. 393, def. 8).

THEOREM 4 (Gelfand--Raikov). --- For every \(x \neq e\) in G, there exists an irreducible unitary representation \(\pi\) of G such that \(\pi(x)\) is not the identity endomorphism of the space \(\pi\) .

Let \(x \in G\) be such that \(\pi(x)\) is the identity for every \(\pi \in \widehat{G}\) . It follows that \(\varphi(x) = \varphi(e) = 1\) for every extreme point \(\varphi\) of \(\mathrm{Pos}_{1}(\mathrm{G})\) (prop. 11 of V, p. 444), and therefore for every \(\varphi \in \mathrm{Pos}_{1}(\mathrm{G})\) by prop. 14, since the linear form \(\varphi \mapsto \varphi(e)\) is continuous on \(\mathrm{Pos}_{1}(\mathrm{G})\) endowed with the weak topology (th. 3). But if \(x \neq e\) , there exists \(f \in L^{2}(\mathrm{G})\) of norm 1 such that \(\langle f | \boldsymbol{\gamma}_{\mathrm{G}}(x)f \rangle = 0\) (V, p. 406, lemma 3). Since the left-hand side of this equality is of the form \(\varphi(x)\) , where \(\varphi \in \mathrm{Pos}_{1}(\mathrm{G})\) , this is a contradiction.

If G is not locally compact, it is not always true that the irreducible unitary representations of G separate the points of G.

COROLLARY. --- For all elements \(g \neq h\) in G, there exists an irreducible representation \(\pi \in \widehat{G}\) and a matrix coefficient f of \(\pi\) such that \(f(g) \neq f(h)\) . In particular, the subalgebra \(\Upsilon(\mathrm{G})\) of \(\mathcal{C}_{b}(\mathrm{G})\) separates the points of G.

Let \(g \neq h\) be in G. There exists an irreducible representation \(\pi \in \widehat{G}\) such that \(\pi(g) \neq \pi(h)\) (th. 4), so there exist x and y in the space of \(\pi\) such that \(\langle x | \pi(g)y \rangle \neq \langle x | \pi(h)y \rangle\) .

\subsection{Functions of positive type on a locally compact commutative group}

In this number, G is a locally compact commutative group and \(\mu\) denotes a Haar measure on G. We denote by \(\widehat{G}\) the dual group of G and by \(\widehat{\mu}\) the dual Haar measure of \(\mu\) (def. 4 of II, p. 214). We denote by F the Fourier transform on the Banach space \(\mathcal{M}^{1}(G)\) of bounded measures on G.

Since \(\mathcal{C}_0(\mathrm{G})\) is the closure of \(\mathcal{K}(\mathrm{G})\) in \(\mathcal{C}_b(\mathrm{G})\) , the Banach space \(\mathcal{M}^1 (\mathrm{G})\) , dual of \(\mathcal{K}(\mathrm{G})\) endowed with the topology of \(\mathcal{C}_b(\mathrm{G})\) , is identified with the dual of \(\mathcal{C}_0(\mathrm{G})\) (cf. INT, III, p. 56, § 1, no. 8). We endow \(\mathcal{M}^1 (\mathrm{G})\) with the weak topology associated with this duality.

Lemma 9. --- The Fourier transform is a continuous mapping from \(\mathcal{M}^{1}(\mathrm{G})\) equipped with the weak topology to \(\mathcal{C}_{b}(\widehat{\mathrm{G}})\) endowed with the topology induced by the weak topology \(\sigma(\mathrm{L}^{\infty}(\widehat{\mathrm{G}}),\mathrm{L}^{1}(\widehat{\mathrm{G}}))\) .

Let \(\mathfrak{F}\) be a filter on \(\mathcal{M}^1 (\mathrm{G})\) weakly converging to a bounded measure \(\nu\) on G. For every \(\varphi \in \mathrm{L}^1 (\widehat{\mathrm{G}})\) , we have \(\mathcal{F}(\varphi)\in \mathcal{C}_0(\mathrm{G})\) (prop. 4 of II, p. 209), hence

\[
\begin{array}{r l} & {\int_ {\widehat {\mathrm{G}}} \varphi \mathcal {F} (\nu) d \widehat {\mu} = \int_ {\mathrm{G}} \mathcal {F} (\varphi) d \nu} \\ & {\qquad = \lim _ {\eta , \mathfrak {F}} \int_ {\mathrm{G}} \mathcal {F} (\varphi) d \eta = \lim _ {\eta , \mathfrak {F}} \int_ {\widehat {\mathrm{G}}} \varphi \mathcal {F} (\eta) d \widehat {\mu},} \end{array}
\]

by the transposition property of the Fourier transform (prop. 13 of II, p. 221). The lemma follows.

THEOREM 5 (Bochner). --- A continuous function \(\varphi\) on \(\widehat{\mathbf{G}}\) belongs to \(\operatorname{Pos}(\widehat{\mathbf{G}})\) if and only if there exists a bounded positive measure \(\nu\) on \(\mathbf{G}\) such that \(\varphi = \mathcal{F}(\nu)\) .

Let \(\nu \in \mathcal{M}^{1}(\mathrm{G})\) be a positive measure. Its Fourier transform is continuous. For every finite family \((x_{i})_{i\in \mathrm{I}}\) in \(\widehat{\mathbf{G}}\) and every finite family \((t_i)_{i\in \mathrm{I}}\) of complex numbers, it follows that

\[
\begin{array}{r} \sum_ {i \in \mathrm{I}} \sum_ {j \in \mathrm{I}} \bar {t} _ {i} t _ {j} \mathcal {F} (\nu) (x _ {i} ^ {- 1} x _ {j}) = \sum_ {i \in \mathrm{I}} \sum_ {j \in \mathrm{I}} \bar {t} _ {i} t _ {j} \int_ {\mathrm{G}} x _ {i} \overline {{x}} _ {j} d \nu \\ = \int_ {\mathrm{G}} \left| \sum_ {i \in \mathrm{I}} \bar {t} _ {i} x _ {i} \right| ^ {2} d \nu \geqslant 0, \end{array}
\]

since \(\nu\) is a positive measure. The Fourier transform of \(\nu\) is therefore a function of positive type on \(\hat{G}\) (th. 1 of V, p. 432, (ii)).

Let us prove the converse. Endow the set \(\operatorname{Pos}_{\leqslant1}(\widehat{\mathbf{G}})\) with the weak topology, induced as before by the weak topology \(\sigma(\mathrm{L}^{\infty}(\widehat{\mathrm{G}}),\mathrm{L}^{1}(\widehat{\mathrm{G}}))\) . The set \(\operatorname{Pos}_{\leqslant1}(\widehat{\mathbf{G}})\) is compact and convex (cor. of prop. 13 of V, p. 448). By prop. 14 of V, p. 450, prop. 11 of V, p. 444, and cor. 7 of V, p. 390, the extreme points of \(\operatorname{Pos}_{\leqslant1}(\widehat{\mathbf{G}})\) are the zero function and the elements of \(\widehat{G}\) .

Let N be the set of positive bounded measures of mass \(\leqslant 1\) on G; it is compact in \(\mathcal{M}^{1}(G)\) endowed with the weak topology (EVT, III, p. 17, cor. 3). By Lemma 9 and the first part of the proof, the Fourier transform on G defines, by passing to subspaces, a continuous map from \(\mathcal{N}\) to \(\mathrm{Pos}_{\leqslant 1}(\widehat{\mathrm{G}})\) ; by homogeneity, it suffices to show that this map is surjective. The image \(\mathcal{F}(\mathcal{N})\) of \(\mathcal{N}\) by the Fourier transform is convex and compact; it contains the zero function and the elements of \(\widehat{\mathrm{G}}\) (indeed, these are of the form \(\mathrm{ev}_x\colon \chi \mapsto \chi (x)\) for an element \(x\) of G by th. 2 of II, p. 220, and we have \(\mathrm{ev}_x = \mathcal{F}(\varepsilon_{x^{-1}})\) ). The set \(\mathcal{F}(\mathcal{N})\) therefore contains the extreme points of \(\mathrm{Pos}_{\leqslant 1}(\mathrm{G})\) , whence \(\mathcal{F}(\mathcal{N}) = \mathrm{Pos}_{\leqslant 1}(\mathrm{G})\) by the Krein-Milman theorem (EVT, II, p. 59, th. 1). This concludes the proof.

When \(G = R^{k}\) for an integer \(k \in N\) , this theorem corresponds to prop. 11 of INT, IX, p. 94, § 6, no. 12.

\section{REPRESENTATIONS OF COMPACT GROUPS}

In this paragraph, all topological vector spaces considered are complex, unless explicitly stated otherwise.

We fix a compact topological group G, whose identity element we denote by e. The group G is unimodular (INT, VII, p. 20, § 1, n° 3, cor. of prop. 3). We denote by μ the normalized Haar measure on G (that is, the one such that μ(G) = 1), and for \(1 \leqslant p \leqslant +\infty\) , we denote by \(\mathcal{L}^{p}(G)\) (resp. \(\mathrm{L}^{p}(\mathrm{G})\) ) the space \(\mathcal{L}_{\mathbf{C}}^{p}(\mathrm{G}, \mu)\) (resp. the space \(\mathrm{L}_{\mathbf{C}}^{p}(\mathrm{G}, \mu)\) ). Convolutions will always be taken with respect to the measure μ. Let \(p \in [1, +\infty]\) . We identify \(\mathcal{C}(\mathrm{G})\) with a subspace of \(\mathrm{L}^{p}(\mathrm{G})\) , which is permissible since the support of μ is equal to G.

For every irreducible unitary representation \(\pi \in \widehat{\mathbf{G}}\) , let \(\mathrm{E}_{\pi}\) be the space of \(\pi\) .

\subsection{Semi-simplicity of finite-dimensional representations}

We recall (INT, VII, p. 71, § 3, n° 1, lemma 1) that for every continuous representation \(\varrho\) of G in a Hilbert space E, there exists a nondegenerate positive Hermitian form \(q\) on E such that the topological vector space structure of E defined by q is identical to the initial structure of E, and such that \(\varrho\) is a unitary representation of G in the Hilbert space E equipped with the inner product q.

PROPOSITION 1. --- Let \(\varrho\) be a finite-dimensional linear representation of G in a separated topological vector space E. There exists an inner product \(q\) on E such that \(\varrho\) is a unitary representation in the Hilbert space E equipped with \(q\) . In particular, \(\varrho\) is semisimple.

Since E is finite-dimensional, there exists a Hilbert space structure on E. The result follows by applying the preceding remark and Corollary 2 of V, p. 392.

Remark. --- We shall see later (Cor. 4 of V, p. 466) that every unitary representation of G is semi-simple.

COROLLARY. --- Let \(\varrho_{1}\) and \(\varrho_{2}\) be finite-dimensional representations of G. The representations \(\varrho_{1}\) and \(\varrho_{2}\) are isomorphic if and only if their characters are equal.

This follows from the proposition and Corollary 3 of V, p. 392.

\subsection{Irreducible representations}

Lemma 1. --- Every irreducible unitary representation of G is square-integrable.

Indeed, the matrix coefficients of an irreducible unitary representation are continuous and bounded, and are therefore square-integrable on G, since G is compact.

PROPOSITION 2. --- Let π be an irreducible unitary representation of G in a Hilbert space E. The dimension of E is finite and equal to the formal degree of π.

The representation \(\pi\) being square-integrable (Lemma 1), its formal degree \(c_{\mu}(\pi)\) relative to \(\mu\) is defined (def. 4 of V, p. 423); it is a strictly positive real number. Let \((e_i)_{i \in I}\) be a finite orthonormal family in E. Let \(x\) be an element of E with norm 1 (such exists since E is nonzero). For \(i \in I\) , we have

\[
c _ {\mu} (\pi) \int_ {\mathrm{G}} | \langle e _ {i} | \pi (g) x \rangle | ^ {2} d \mu (g) = 1
\]

(prop. 8 of V, p. 424). Summing this formula over \(i \in I\) , we obtain

\[
\begin{array}{r l} & {\mathrm{Card(I)} = c _ {\mu} (\pi) \int_ {\mathrm{G}} \sum_ {i \in \mathrm{I}} | \langle e _ {i} | \pi (g) x \rangle | ^ {2} d \mu (g)} \\ & {\qquad \leqslant c _ {\mu} (\pi) \int_ {\mathrm{G}} \| \pi (g) x \| ^ {2} d \mu (g) = c _ {\mu} (\pi)} \end{array}
\]

by Bessel's inequality (EVT, V, p. 21, prop. 4). Hence the cardinality of I is bounded above by \(c_{\mu}(\pi)\) . This implies that the dimension of E is finite.

One can then apply the preceding result to an orthonormal basis \((e_{i})_{i\in\mathrm{I}}\) of E. We obtain, by EVT, V, p. 22, prop. 5,

\[
\begin{array}{r} \dim (\mathrm{E}) = c _ {\mu} (\pi) \int_ {\mathrm{G}} \sum_ {i \in \mathrm{I}} | \langle e _ {i} | \pi (g) x \rangle | ^ {2} d \mu (g) \\ = c _ {\mu} (\pi) \int_ {\mathrm{G}} \| \pi (g) x \| ^ {2} d \mu (g) = c _ {\mu} (\pi), \end{array}
\]

which concludes the proof.

In particular, it is therefore permissible to speak of the character \(\chi_{\pi}\) of an irreducible unitary representation \(\pi\) of G. It is a continuous function on G.

COROLLARY 1. --- Let \(\pi\) be an irreducible unitary representation of G in a Hilbert space E. The character of \(\pi\) is a Hermitian element of the involutive algebra \(\mathrm{L}^{1}(\mathrm{G})\) .

Let \(x \in G\) . Since \(G\) is unimodular, we have \(\chi_{\pi}^{*}(x) = \overline{\operatorname{Tr}(\pi(x^{-1}))}\) , whence \(\chi_{\pi}^{*}(x) = \operatorname{Tr}(\pi(x))\) since \(\pi\) is a unitary representation.

COROLLARY 2. --- a) Let \(\pi\) be an irreducible unitary representation of G in a Hilbert space E. We have

\[
\int_ {\mathrm{G}} \overline {{\langle x \mid \pi (g) x ^ {\prime} \rangle}} \langle y \mid \pi (g) y ^ {\prime} \rangle d \mu (g) = \frac {1}{\dim (\mathrm{E})} \overline {{\langle x \mid y \rangle}} \langle x ^ {\prime} \mid y ^ {\prime} \rangle
\]

for all \((x,y,x',y')\in \mathrm{E}^4\) .

\begin{enumerate}
\def\labelenumi{\alph{enumi})}
\setcounter{enumi}{1}
\tightlist
\item
  Let \(\pi_{1}\) (resp. \(\pi_{2}\) ) be an irreducible unitary representation of G in a Hilbert space \(E_{1}\) (resp. \(E_{2}\) ). If \(\pi_{1}\) and \(\pi_{2}\) are not isomorphic, one has
\end{enumerate}

\[
\int_ {G} \overline {{\langle x \mid \pi_ {1} (g) x ^ {\prime} \rangle}} \langle y \mid \pi_ {2} (g) y ^ {\prime} \rangle d \mu (g) = 0
\]

for all \((x,x^{\prime},y,y^{\prime})\in \mathrm{E}_{1}^{2}\times \mathrm{E}_{2}^{2}\)

This follows immediately from prop. 8 of V, p. 424 and prop. 9 of V, p. 424, taking into account lemma 1 and the formula \(c_{\mu}(\pi) = \dim(\mathrm{E})\) (prop. 2).

COROLLARY 3. --- Let \(\pi_1\) and \(\pi_2\) be irreducible representations of G. In the Hilbert space \(\mathrm{L}^{2}(\mathrm{G})\) , one has \(\langle \chi_{\pi_1} | \chi_{\pi_2} \rangle = 1\) if \(\pi_1\) and \(\pi_2\) are isomorphic and \(\langle \chi_{\pi_1} | \chi_{\pi_2} \rangle = 0\) otherwise. In other words, the family of characters of the classes \(\pi \in \widehat{\mathbf{G}}\) is an orthonormal family in \(\mathrm{L}^{2}(\mathrm{G})\) .

This follows from proposition 1 of V, p. 457 and corollary 2, noting that for every orthonormal basis \((e_{i})_{i\in\mathrm{I}}\) of the space of a finite-dimensional unitary representation \(\pi\) of G, and for every \(g\in G\) , we have the formula

\[
\chi_ {\pi} (g) = \sum_ {i \in I} \langle e _ {i} | \pi (g) e _ {i} \rangle .
\]

COROLLARY 4. --- a) Let \(\varrho\) be a continuous finite-dimensional representation of G. One has

\[
\chi_ {\varrho} = \sum_ {\pi \in \widehat {G}} \dim (\operatorname{Hom} _ {G} (\pi , \varrho)) \chi_ {\pi} = \sum_ {\pi \in \widehat {G}} \dim (\operatorname{Hom} _ {G} (\varrho , \pi)) \chi_ {\pi}.
\]

\begin{enumerate}
\def\labelenumi{\alph{enumi})}
\setcounter{enumi}{1}
\tightlist
\item
  Let \(\varrho_{1}\) and \(\varrho_{2}\) be continuous finite-dimensional representations of G. One has
\end{enumerate}

\[
\left\langle \chi_ {\varrho_ {1}} \mid \chi_ {\varrho_ {2}} \right\rangle = \dim (\operatorname{Hom} _ {\mathrm{G}} (\varrho_ {1}, \varrho_ {2})) = \dim (\operatorname{Hom} _ {\mathrm{G}} (\varrho_ {2}, \varrho_ {1})).
\]

\begin{enumerate}
\def\labelenumi{\alph{enumi})}
\setcounter{enumi}{2}
\tightlist
\item
  A continuous finite-dimensional representation \(\varrho\) of G is irreducible if and only if \(\| \chi_{\varrho}\|^{2} = 1\) .
\end{enumerate}

Since \(\varrho\) is semi-simple, it is isomorphic to the direct sum of its \(\pi\)-isotypic components \(\mathrm{M}_{\pi}(\varrho)\) for \(\pi \in \widehat{\mathrm{G}}\) . Let \(m_{\pi}(\varrho)\) denote the multiplicity of \(\pi\) in \(\varrho\) (def. 10 of V, p. 398); we then have

\[
\chi_ {\varrho} = \sum_ {\pi \in \widehat {G}} m _ {\pi} (\varrho) \chi_ {\pi}.
\]

The assertion a) therefore follows from the fact that

\[
m _ {\pi} (\varrho) = \dim \operatorname{Hom} _ {\mathrm{G}} (\pi , \varrho) = \dim \operatorname{Hom} _ {\mathrm{G}} (\varrho , \pi)
\]

(formula (1) of V, p. 377 and cor. 2 of V, p. 387).

By bilinearity and orthonormality of characters, the assertion \(a\) ) implies that

\[
\langle \chi_ {\varrho_ {1}} \mid \chi_ {\varrho_ {2}} \rangle = \sum_ {\pi \in \widehat {\mathrm{G}}} \dim (\operatorname{Hom} _ {\mathrm{G}} (\pi , \varrho_ {1})) \dim (\operatorname{Hom} _ {\mathrm{G}} (\pi , \varrho_ {2})),
\]

and on the other hand we have canonical isomorphisms

\[
\begin{array}{c} \operatorname{Hom} _ {\mathrm{G}} (\varrho _ {1}, \varrho _ {2}) \to \bigoplus_ {(\pi _ {1}, \pi _ {2}) \in \widehat {\mathrm{G}} \times \widehat {\mathrm{G}}} \operatorname{Hom} _ {\mathrm{G}} (\mathrm{M} _ {\pi _ {1}} (\varrho _ {1}), \mathrm{M} _ {\pi _ {2}} (\varrho _ {2})) \\ \qquad \qquad \qquad \qquad \qquad \qquad \qquad \qquad \qquad \qquad \qquad \qquad \qquad \qquad \qquad \qquad \qquad \qquad \qquad \qquad \qquad \qquad \qquad \qquad \qquad \qquad \qquad \qquad \qquad \qquad \qquad \qquad \qquad \qquad \qquad \to \bigoplus_ {\pi \in \widehat {\mathrm{G}}} \operatorname{Hom} _ {\mathrm{G}} (\mathrm{M} _ {\pi} (\varrho _ {1}), \mathrm{M} _ {\pi} (\varrho _ {2})) \end{array}
\]

(formula (1) of V, p. 377), whence it follows that

\[
\dim (\operatorname{Hom} _ {\mathrm{G}} (\varrho_ {1}, \varrho_ {2})) = \sum_ {\pi \in \widehat {\mathrm{G}}} m _ {\pi} (\varrho_ {1}) m _ {\pi} (\varrho_ {2}),
\]

whence the assertion \(b\) ). The assertion \(c\) ) also follows from \(a\) ) and Schur's lemma (prop. 6 of V, p. 386).

COROLLARY 5. --- Let \(\pi_1\) and \(\pi_2\) be irreducible unitary representations of G. Then \(\pi_1(\overline{\chi}_{\pi_2}) = 0\) if \(\pi_1\) is not isomorphic to \(\pi_2\) , and \(\pi_1(\overline{\chi}_{\pi_1})\) is multiplication by \(1/\dim(\pi_1)\) .

The character of \(\pi_2\) is a continuous central function on G, hence the linear map \(u = \pi_1(\overline{\chi}_{\pi_2})\) is defined and belongs to the space \(\mathrm{Hom}_{\mathrm{G}}(\pi_1,\pi_1)\) . It is a homothety by Schur's lemma (prop. 6 of V, p. 386) whose trace is

\[
\begin{array}{c} \operatorname{Tr} (u) = \operatorname{Tr} \Bigl (\int_ {\mathrm{G}} \overline {{\chi_ {\pi_ {2}} (g)}}   \pi_ {1} (g) d \mu (g) \Bigr) \\ = \int_ {\mathrm{G}} \overline {{\chi_ {\pi_ {2}} (g)}} \chi_ {\pi_ {1}} (g) d \mu (g). \end{array}
\]

By the orthogonality relations, the trace of \(u\) is therefore zero if \(\pi_1\) is not isomorphic to \(\pi_2\) and equals 1 otherwise. The assertion follows.

Remark. --- When G is finite, the Schur orthogonality relation (resp. the character orthogonality formula) coincides with that of A, VIII, p. 399 (resp. that of A, VIII, p. 400, prop. 4). On the other hand, the "second orthogonality relation" of characters of finite groups (A, VIII, p. 402, formula (32)) has no exact analogue when G is compact.

For G finite, the particular case of the second orthogonality formula corresponding to the conjugacy class of e is the formula

\[
\frac {1}{\mathrm{CardG}} \sum_ {\pi \in \widehat {\mathrm{G}}} \dim (\pi) \chi _ {\pi} (g) = \left\{ \begin{array}{l l} 1 & \text {if} g = e \\ 0 & \text {otherwise}, \end{array} \right.
\]

which is also interpreted as the computation of the character of the regular representation \(\gamma_{G}\) of G, and which is equivalent to the formula \(\mathrm{Tr}(\boldsymbol{\gamma}_{\mathrm{G}}(f)) = f(e)\) for every function f from G to C.

Let G again be an arbitrary compact group. For every function \(f \in \mathcal{C}(\mathrm{G})\) , the endomorphism \(\boldsymbol{\gamma}_{\mathrm{G}}(f)\) of \(\mathrm{L}^{2}(\mathrm{G})\) coincides with the endomorphism \(\varphi \mapsto f * \varphi\) and it has finite trace (lemma 4 of V, p. 407 and corollary 2 of III, p. 33). By loc. cit. and th. 2 of IV, p. 177, we also have

\[
\operatorname{Tr} \left(\boldsymbol {\gamma} _ {\mathrm{G}} (f)\right) = \int_ {\mathrm{G}} f \left(x ^ {- 1} x\right) d \mu (x) = f (e).
\]

PROPOSITION 3. --- Suppose that \(G = G_{1} \times G_{2}\) where \(G_{1}\) and \(G_{2}\) are compact groups. The map b from \(\widehat{G}_{1} \times \widehat{G}_{2}\) to \(\widehat{G}\) that sends \((\pi_{1}, \pi_{2})\) to the class of \(\pi_{1} \boxtimes \pi_{2}\) is a bijection.

By cor. 6 of V, p. 389, the map b is well defined.

Let \(\pi_1\) and \(\pi_2\) be elements of \(\widehat{\mathbf{G}}\) . Let \(\pi \in \widehat{\mathbf{G}}\) . The \(\pi\)-isotypic component of the restriction \(\varpi\) of \(\pi_1 \boxtimes \pi_2\) to \(\mathrm{G}_1 \times \{e\}\) is equal to \(\varpi\) if \(\pi_1 = \pi\) and is zero otherwise (lemma 8 of V, p. 384). Thus, the unitary representation \(\pi_1 \boxtimes \pi_2\) determines \(\pi_1\) up to isomorphism; likewise, it determines \(\pi_2\) , which proves that \(b\) is injective.

Let us prove that \(b\) is surjective. Let \(\pi\) be an irreducible unitary representation of \(G_1 \times G_2\) in a Hilbert space E. It is finite-dimensional (prop. 2). Let \(\pi_1\) be an irreducible unitary representation of \(G_1\) in a Hilbert space \(E_1\) such that the restriction of \(\pi\) to \(G_1 \times \{e\}\) contains a subrepresentation isomorphic to \(\pi_1\) (lemma 3 of V, p. 379). Let \(F = \text{Hom}_{G_1}(\pi_1, \pi)\) . This is a nonzero finite-dimensional vector space. For \(h \in G_2\) and \(u \in F\) , let \(\varrho(h)(u)\) be the linear map from \(E_1\) to E defined by \(x \mapsto \pi(e, h)(u(x))\) . Since \(G_1 \times \{e\}\) and \(\{e\} \times G_2\) commute in G, the map \(\varrho(h)(u)\) belongs to the space F. The map \(\varrho\) is a linear representation of \(G_2\) in the space F; it is continuous. Since F is nonzero, there exists an irreducible subrepresentation \(\pi_2\) of \(\varrho\) (lemma 3 of V, p. 379). Let \(E_2\) be the space of \(\pi_2\) . The linear map \(v: E_1 \otimes E_2 \to E\) such that \(x \otimes u \mapsto u(x)\) for \(x \in E_1\) and \(u \in E_2\) is then a nonzero G-morphism from \(\pi_1 \boxtimes \pi_2\) to \(\pi\); since the representations \(\pi\) and \(\pi_1 \boxtimes \pi_2\) are irreducible and finite-dimensional, the morphism \(v\) is an isomorphism (Lemma 2 of V, p. 378).

This proposition is to be compared with Theorem 1 of A, VIII, p. 208.

\subsection{The Peter-Weyl theorem}

We recall that we denote by \(\Theta(\mathrm{G})\) the space of matrix coefficients of finite-dimensional unitary representations of G. The space \(\Theta(\mathrm{G})\) is a unital subalgebra of \(\mathcal{C}(\mathrm{G})\) stable under complex conjugation (prop. 5 of V, p. 386).

For \(\pi \in \widehat{\mathrm{G}}\) , let \(\varrho_{\mathrm{G}}(\pi)\) be the subspace of \(\mathcal{C}(\mathrm{G})\) generated by the matrix coefficients of \(\pi\) . We identify it with a subspace of \(\mathrm{L}^2 (\mathrm{G})\) .

The space \(\Theta(G)\) coincides with the sum of the spaces \(\varrho_{G}(\pi)\) for \(\pi\in\widehat{G}\) . Indeed, every element of \(\Theta(G)\) is a sum of matrix coefficients of irreducible representations of G, since the finite-dimensional representations of G are semisimple by Prop. 1 of V, p. 457). Moreover, this sum is direct since, by Cor. 2 of V, p. 458, the spaces \(\varrho_{G}(\pi)\) are pairwise orthogonal.

Proposition 4. — The space \(\Theta(\mathrm{G})\) is dense in \(\mathcal{C}(\mathrm{G})\) and in \(\mathrm{L}^2(\mathrm{G})\) .

The unital subalgebra \(\Theta(\mathrm{G})\) of \(\mathcal{C}(\mathrm{G})\) is stable under conjugation. It coincides with the subalgebra \(\Upsilon(\mathrm{G})\) by prop. 2 of V, p. 457, hence it separates the points of G (cor. of th. 4 of V, p. 454). Consequently, it is dense in \(\mathcal{C}(\mathrm{G})\) by TG, X, p. 40, cor. 2, and a fortiori, it is dense in the space \(\mathrm{L}^2(\mathrm{G})\) (cf. INT, IV, p. 155, §4, no. 7, prop. 13).

Recall that the biregular unitary representation of G is the representation \(\varrho_{G}\) of G × G into \(\mathrm{L}^{2}(\mathrm{G})\) such that \((g_{1}, g_{2}) \mapsto \gamma_{\mathrm{G}}(g_{1}) \delta_{\mathrm{G}}(g_{2})\) .

PROPOSITION 5. --- Let \(\pi \in \widehat{\mathbf{G}}\) . The space \(\varrho_{\mathrm{G}}(\pi)\) is a subrepresentation of \(\varrho_{\mathrm{G}}\) which is isomorphic to \(\overline{\pi} \boxtimes \pi\) . In particular, it is an irreducible representation of \(\mathrm{G} \times \mathrm{G}\) .

By prop. 7 of V, p. 422 (applied with \(Z = \{e\}\) ), the space \(\varrho_{G}(\pi)\) is a subrepresentation of \(\varrho_{G}\) and the linear map from \(\overline{\mathrm{E}}_{\pi} \otimes \mathrm{E}_{\pi}\) to \(L^2(G)\) that sends \(x \otimes y\) to the matrix coefficient \(g \mapsto \langle x | \pi(g)y \rangle\) is a \((G \times G)\)-isomorphism from \(\overline{\pi} \boxtimes \pi\) onto \(\varrho_{G}(\pi)\) . The subrepresentation \(\varrho_{G}(\pi)\) is therefore irreducible (prop. 3 of V, p. 461).

THEOREM 1 (Peter--Weyl). --- Let G be a compact topological group. The biregular representation \(\varrho_{\mathrm{G}}\) of G is the Hilbert sum of the subrepresentations \(\varrho_{\mathrm{G}}(\pi)\) for \(\pi \in \widehat{\mathrm{G}}\) .

The spaces \(\varrho_{\mathrm{G}}(\pi)\) are pairwise orthogonal; these are irreducible subrepresentations of the biregular representation of G (prop. 5) which are pairwise non-isomorphic (prop. 3 of V, p. 461).

The theorem then follows from the fact that the sum \(\Theta(\mathrm{G})\) of the spaces \(\varrho_{\mathrm{G}}(\pi)\) for \(\pi\) running through \(\widehat{\mathrm{G}}\) is dense in \(\mathrm{L}^2(\mathrm{G})\) (prop. 4).

COROLLARY 1. --- For every \(\pi\in\widehat{G}\) , the subrepresentation \(\varrho_{\mathrm{G}}(\pi)\) is the \((\overline{\pi}\boxtimes\pi)\)-isotypic component of \(\varrho_{\mathrm{G}}\) .

Let F be the \((\overline{\pi} \boxtimes \pi)\)-isotypic component of \(\varrho_{G}\) . The space F contains \(\varrho_{\mathrm{G}}(\pi)\) (prop. 5). Moreover, for every \(\tau \in \widehat{G}\) different from \(\pi\) , the intersection of F and \(\varrho_{\mathrm{G}}(\tau)\) is zero (prop. 3 of V, p. 461). From this we deduce that \(\mathrm{F} = \varrho_{\mathrm{G}}(\pi)\) by applying the theorem.

COROLLARY 2. --- Let \(\pi_1\) and \(\pi_2\) be non-isomorphic irreducible representations of G. The \((\overline{\pi}_1 \boxtimes \pi_2)\)-isotypic component of \(\varrho_G\) is zero.

COROLLARY 3. --- The right (resp. left) regular representation of G is isomorphic to the Hilbert sum

\[
\bigoplus_ {\pi \in \widehat {\mathrm{G}}} \pi^ {\dim (\pi)}.
\]

By the theorem and proposition 5, the restriction of the biregular representation of G to the subgroup \(\mathrm{H} = \{e\} \times \mathrm{G}\) is isomorphic to the Hilbert sum of the restrictions to H of the representations \(\overline{\pi} \boxtimes \pi\) for \(\pi \in \widehat{\mathrm{G}}\) . These are isomorphic to the direct sum of \(\dim(\pi)\) copies of \(\pi\) (lemma 8 of V, p. 384). This implies the result for the right regular representation, and the case of the left regular representation is similar.

COROLLARY 4. --- Let \(\pi\in\widehat{G}\) . The \(\pi\)-isotypic component of \(\delta_{G}\) (resp. the component \(\overline{\pi}\)-isotypic of \(\gamma_{G}\) ) is equal to \(\varrho_{G}(\pi)\) .

The argument is similar to that of the preceding corollary, by considering \(\varrho_{\mathrm{G}}(\pi)\) as a subrepresentation of \(\delta_{\mathrm{G}}\) (resp. of \(\gamma_{\mathrm{G}}\) ).

Remarks. --- 1) If G is finite, these statements correspond to the results of A, VIII, p. 398, remark.

\begin{enumerate}
\def\labelenumi{\arabic{enumi})}
\setcounter{enumi}{1}
\item
  Suppose that G is commutative. The set \(\widehat{G}\) identifies with the dual group of G (V, p. 393, remark). For every \(\chi \in \widehat{G}\) , the space \(\varrho_{\mathrm{G}}(\chi)\) is the 1-dimensional subspace of \(\mathrm{L}^2 (\mathrm{G})\) generated by \(\chi\) . Theorem 1 for G is therefore equivalent to the corollary of Theorem 1 of II, p. 215.
\item
  If there exists an integer \(n \geqslant 0\) such that G is a compact subgroup of \(\mathbf{GL}(\mathbf{C}^{n})\) , the identity representation of G in \(\mathbf{GL}(\mathbf{C}^{n})\) is enough to separate the points of G, and one can then prove Proposition 4 directly, followed by the Peter–Weyl theorem, without invoking the Gelfand–Raikov theorem.
\item
  The Peter–Weyl theorem implies in particular that the natural continuous homomorphism from G into \(\prod_{\pi\in\widehat{G}}\mathbf{U}(\mathrm{E}_{\pi})\) is injective.
\end{enumerate}

\subsection{Matrix coefficients and G-finite functions}

PROPOSITION 6. --- The following subspaces of \(\mathrm{L}^{2}(\mathrm{G})\) are equal:

\begin{enumerate}
\def\labelenumi{\alph{enumi})}
\item
  The space \(\Theta (\mathbf{G})\)
\item
  The algebraic direct sum of the subspaces \(\varrho_{\mathrm{G}}(\pi)\) of \(\mathrm{L}^2 (\mathrm{G})\) .
\item
  The space of G-finite vectors (cf. V, p. 376) of \(\gamma_{\mathrm{G}}\) ;
\item
  The space of G-finite vectors of \(\delta_{G}\) ;
\item
  The space of \((\mathbf{G} \times \mathbf{G})\)-finite vectors of \(\varrho_{\mathrm{G}}\) .
\end{enumerate}

In particular, every G-finite vector of \(\gamma_{\mathrm{G}}\) , \(\delta_{G}\) or \(\varrho_{\mathrm{G}}\) belongs to \(\mathcal{C}(\mathrm{G})\) .

Denote by \(F_{a}\) (resp. \(F_{b}\) , \(F_{c}\) , \(F_{d}\) , \(F_{e}\) ) the space defined by condition a) (resp. b), c), d), e)). We have already noted that \(F_{a} = F_{b}\) .

One has \(F_{b} \subset F_{c}\) since \(\varrho_{\mathrm{G}}(\pi)\) is a finite-dimensional subrepresentation of \(\gamma_{G}\) for all \(\pi \in \widehat{G}\) . Conversely, let f be a G-finite vector of \(\gamma_{G}\) . The subspace \(E_{f}\) generated by the functions \(\gamma_{\mathrm{G}}(g)f\) for \(g \in G\) is a finite-dimensional subrepresentation of \(\gamma_{G}\) . It is equal to the direct sum of its \(\pi\)-isotypic components for \(\pi \in \widehat{G}\) (prop. 1 of V, p. 457). According to cor. 4 of V, p. 463, this implies that \(E_{f}\) is contained in the sum of the spaces \(\varrho_{\mathrm{G}}(\pi)\) , hence \(F_{c} \subset F_{b}\) .

By an analogous argument, we obtain \(F_{b} = F_{d}\) , and since \(\varrho_{\mathrm{G}}(\pi)\) is a subrepresentation of \(\varrho_{G}\) , one establishes in the same way that \(F_{b} = F_{e}\) .

COROLLARY. --- Let \(E \subset L^{2}(G)\) be a finite-dimensional vector subspace defining a subrepresentation of \(\gamma_{G}\) (resp. of \(\delta_{G}\) , of \(\varrho_{G}\) ). Then E is contained in \(\mathcal{C}(G)\) .

Indeed, every element of E is a G-finite vector of the representation \(\gamma_{G}\) .

\subsection{Representations in a separated quasi-complete space}

PROPOSITION 7. --- Let \(\varrho\) be a continuous representation of G in a separated, locally convex, and quasi-complete space E. The sum of the finite-dimensional subrepresentations of E is dense in E.

Let \(x \in \mathrm{E}\) and let U be an open neighborhood of \(x\) . The set of measures \(\nu \in \mathcal{M}(\mathrm{G})\) such that \(\varrho(\nu)x \in \mathrm{U}\) is open in \(\mathcal{M}(\mathrm{G})\) for the topology of compact convergence (cf. no. 2 of V, p. 400). It contains \(\varepsilon_e\) ; hence it contains a measure of the form \(\nu = f_1 \cdot \mu\) where \(f_1 \in \mathcal{C}(\mathrm{G})\) (INT, VIII, p. 171, § 4, n° 7, prop. 19) and, consequently, it contains a measure of the form \(f_2 \cdot \mu\) where \(f_2\) is a G-finite function (prop. 6 of V, p. 464 and prop. 4 of V, p. 462).

The subrepresentation F of \(\gamma_{G}\) generated by \(f_{2}\) is finite-dimensional. Let \(\widetilde{F}\) be the image of the linear map \(f \mapsto \varrho(f)x\) from F into E. The space \(\widetilde{F}\) is finite-dimensional, and it contains the element \(\varrho(f_{2})x\) belonging to U. Since \(\varrho(g)\varrho(f) = \varrho(\gamma_{\mathrm{G}}(g)f)\) for all \(g \in G\) and every \(f \in F\) (formula (1) of V, p. 406), the space \(\widetilde{F}\) is a subrepresentation of \(\varrho\) which intersects U. The proposition is proved.

COROLLARY 1. --- Let π be a continuous irreducible representation of G in a locally convex separated quasi-complete space E. The space E is finite-dimensional.

COROLLARY 2. --- Let \(\varrho\) be a continuous representation of G in a locally convex separated quasi-complete space E, and let \(\pi\) be a continuous irreducible representation of G.

\begin{enumerate}
\def\labelenumi{\alph{enumi})}
\item
  The G-morphism \(p_{\pi} = \dim(\pi)\varrho(\overline{\chi}_{\pi})\) from E into E is a continuous projector of E whose image is the \(\pi\)-isotypic component of \(\varrho\) ;
\item
  If \(\varrho\) is a unitary representation, then the projector \(p_{\pi}\) is the orthogonal projector of E with image \(\mathrm{M}_{\pi}(\varrho)\) .
\end{enumerate}

Let us prove \(a\) ). The linear map \(p_{\pi} = \dim(\pi)\varrho(\overline{\chi}_{\pi})\) is well defined, since E is quasi-complete and G is compact (V, p. 401). It is an element of \(\mathrm{Hom}_{\mathrm{G}}(\varrho,\varrho)\) since the character of \(\pi\) is a continuous central function on G.

Let \(\varpi\) be a finite-dimensional subrepresentation of E and F its space. The map \(p_{\pi}\) induces, by passing to subspaces, the endomorphism \(\dim(\pi)\varpi(\overline{\chi}_{\pi})\) of F. This endomorphism is therefore zero if \(\varpi\) is not isomorphic to \(\pi\) and is the identity map of F otherwise (indeed, this is the case if \(\varpi\) is irreducible by cor. 5 of V, p. 460, and the general case follows from this since \(\varpi\) is semisimple by prop. 1 of V, p. 457).

Finally, since the sum of the finite-dimensional subrepresentations of E is dense in E (prop. 7) and since \(p_{\pi}\) is continuous, we conclude that \(p_{\pi}\) is a projector whose image is the \(\pi\)-isotypic component of \(\varrho\) . If E is a Hilbert space, the projector \(p_{\pi}\) is Hermitian by Cor. 1 of V, p. 458, hence it is an orthogonal projector (Lemma 3, (ii) of I, p. 133).

COROLLARY 3. --- Let \(\varrho\) be a continuous representation of G in a locally convex separated quasi-complete space E. The endomorphism

\[
x \mapsto \int_ {\mathrm{G}} \varrho (g) x d \mu (g)
\]

of E is a projector of E whose image is the space \(E^{G}\) of invariant vectors in E. In particular, if E is finite-dimensional, then

\[
\mathrm{dim} (\mathrm{E} ^ {\mathrm{G}}) = \int_ {\mathrm{G}} \chi_ {\varrho} (g) d \mu (g).
\]

The first assertion is the special case of the preceding corollary when \(\pi\) is the trivial representation of G, of dimension 1. The second follows from it since the dimension of \(E^{G}\) is the trace of the orthogonal projector onto \(E^{G}\) , that is,

\[
\mathrm{dim} (\mathrm{E} ^ {\mathrm{G}}) = \mathrm{Tr} \Bigl (\int_ {\mathrm{G}} \varrho (g) d \mu (g) \Bigr) = \int_ {\mathrm{G}} \chi_ {\varrho} (g) d \mu (g).
\]

In particular, when E is finite-dimensional, there exists a vector \(x \neq 0\) in E such that \(\varrho(g)x = x\) for all \(g \in G\) if and only if

\[
\int_ {\mathrm{G}} \chi_ {\varrho} (g) d \mu (g) \neq 0.
\]

COROLLARY 4. --- Let \(\varrho\) be a unitary representation of G in a Hilbert space E. The space E is the Hilbert sum of the \(\pi\)-isotypic components \(\mathrm{M}_{\pi}(\varrho)\) for \(\pi \in \widehat{\mathrm{G}}\) . In particular, the representation \(\varrho\) is semisimple.

The isotypic components of \(\varrho\) corresponding to non-isomorphic irreducible representations of G are orthogonal (prop. 8 of V, p. 394). Since every finite-dimensional unitary representation of G is semisimple (prop. 1 of V, p. 457), the sum of the isotypic components \(M_{\pi}(\varrho)\) for \(\pi \in \widehat{G}\) is dense in E (prop. 7). The corollary follows from this.